\documentclass[a4paper,12pt]{article}

% Pacchetti utili
\usepackage[italian]{babel}
\usepackage[T1]{fontenc}
\usepackage[utf8]{inputenc}
\usepackage{lmodern}
\usepackage{geometry}
\usepackage{setspace}
\usepackage{graphicx}
\usepackage{titling}
\usepackage{hyperref}
\usepackage{booktabs}
\usepackage{amsmath, amssymb}
\usepackage[table]{xcolor}
\usepackage[most]{tcolorbox}
\usepackage{tikz}
\usepackage{enumitem}

\geometry{margin=2.5cm}
\onehalfspacing

\begin{document}
	
	\begin{titlepage}
		\centering
		
		% Intestazione Università
		{\scshape\LARGE Università di Pisa \par}
		\vspace{0.5cm}
		
		% Titolo corso
		{\Huge\bfseries Economia ed Organizzazione Aziendale \par}
		\vspace{0.5cm}
		{\Large Anno Accademico 2025--2026 \par}
		
		\vfill
		
		% Autore
		{\Large\textbf{Pietro Bacigalupi} \par}
		\vspace{0.5cm}
		
		\vfill
		
		% Data
		{\large \today\par}
		
	\end{titlepage}
	
	\tableofcontents
	\newpage
	
	\section{Introduzione}
	
	\subsection{Cos'è l'economia?}
	
	L'economia può essere definita sinteticamente come segue:
	
	\begin{tcolorbox}[colback=blue!5,colframe=blue!75!black,title=Definizione di Economia]
		\textbf{Scienza che studia come gli individui e le società scelgono di allocare risorse scarse tra usi alternativi per soddisfare bisogni illimitati.}
	\end{tcolorbox}
	
	Si tratta di una \textbf{scienza} in quanto approccia lo studio dei fenomeni attraverso il metodo scientifico, ma è una \textbf{scienza sociale}, poiché il suo oggetto di studio è il comportamento umano. 
	
	\subsection{Limiti del modello tradizionale}
	Il comportamento delle persone è spesso meno razionale di quanto i modelli economici tradizionali tendano a ipotizzare. Storicamente, l'economia ha modellizzato l'individuo come un \textit{Homo Economicus}, un agente perfettamente razionale che massimizza una \textbf{funzione di utilità} matematica. 
	
	Tuttavia, questo approccio presenta delle criticità:
	\begin{itemize}
		\item \textbf{Benessere vs Utilità:} Si tende a sovrapporre la funzione di utilità con il concetto di benessere o felicità.
		\item \textbf{Indicatori Macroeconomici:} A livello aggregato, il PIL viene spesso usato come misura dell'utilità totale di un Paese, producendo effetti distorsivi poiché non tiene conto della distribuzione della ricchezza o della qualità della vita.
	\end{itemize}
	
	\subsection{L'importanza delle intenzioni}
	Per comprendere perché il comportamento umano non risponda sempre a logiche puramente matematiche, si può analizzare un celebre esperimento di economia comportamentale.
	
	\begin{quote}
		\textit{Scenario:} Un soggetto A riceve 10€ e deve decidere quanti di questi cedere a un soggetto B. 
		\begin{enumerate}
			\item Se B accetta l'offerta, la somma viene divisa come proposto da A.
			\item Se B rifiuta, entrambi i soggetti ricevono 0€.
		\end{enumerate}
	\end{quote}
	
	Secondo la razionalità pura, B dovrebbe accettare qualsiasi somma maggiore di zero (anche solo 1€), poiché 1€ è comunque meglio di nulla. In realtà, l'esperimento dimostra che:
	\begin{itemize}
		\item Se A offre 2€ avendo la possibilità di offrirne 5 (dimostrando scarsa equità), la maggior parte dei soggetti B \textbf{rifiuta} l'offerta per "punire" il comportamento egoista di A.
		\item Se la scelta di A è invece vincolata (ad esempio, il sistema gli permetteva di offrire solo 0 o 2€), allora B tende ad \textbf{accettare}.
	\end{itemize}
	
	\textbf{Conclusione:} Gli esseri umani non valutano solo le conseguenze materiali (il guadagno monetario), ma anche le \textbf{intenzioni} e l'equità percepita dietro un'azione. Il benessere derivante da uno scambio non è quindi solo quantitativo, ma profondamente legato alla reciprocità sociale.
	
	Questo significa che il nostro comportamento sociale è fortemente guidato dalle reazioni alle intenzioni che noi attribuiamo alle scelte degli altri.
	% da qui 
	\subsection{L'approccio utilitarista e i suoi limiti}
	Il modello economico classico si basa spesso sulla logica dell' \textbf{utilitarismo}. Secondo questa visione, un'azione (o una politica economica) è valutata come "buona" o "efficace" esclusivamente in base alla sua capacità di aumentare la \textbf{funzione di utilità totale} della società.
	
	Tuttavia, questo approccio presenta un limite etico e analitico fondamentale:
	\begin{itemize}
		\item \textbf{Indifferenza alla distribuzione:} Poiché si guarda solo alla somma totale, il risultato matematico non cambia se l'utilità di un individuo viene azzerata per aumentare drasticamente quella di un altro. 
		\item \textbf{Assenza di equità:} Se ci basiamo solo sull'utilità totale, non siamo in grado di distinguere tra una società dove tutti stanno mediamente bene e una società dove pochi hanno moltissimo e molti nulla, purché la "somma" sia la medesima.
	\end{itemize}
	
	\subsection{Economia come sistema complesso}
	In ambito gestionale e organizzativo, è fondamentale distinguere tra sistemi \textbf{complicati} e sistemi \textbf{complessi}. Spesso questi termini vengono usati come sinonimi, ma identificano realtà profondamente diverse.
	
	\subsubsection{Il sistema complicato}
	Un sistema complicato (come un orologio Patek Philippe, un motore aeronautico o un software avanzato) è caratterizzato da:
	\begin{itemize}
		\item \textbf{Riduzionismo:} Il sistema può essere compreso pienamente smontandolo e analizzando le singole $n$ parti che lo compongono.
		\item \textbf{Linearità:} L'insieme è esattamente la \textbf{somma delle parti}. Non vi sono proprietà emergenti impreviste.
		\item \textbf{Prevedibilità:} È un sistema deterministico. A fronte di una determinata azione (input), l'effetto (output) è noto a priori e costante nel tempo.
		\item \textbf{Gestione:} Viene gestito attraverso il controllo e la pianificazione tecnica.
	\end{itemize}
	
	\subsubsection{Il sistema complesso}
	L'economia e le organizzazioni umane sono invece sistemi complessi. Le loro caratteristiche sono:
	\begin{itemize}
		\item \textbf{Relazioni dinamiche:} Le $n$ parti che compongono il sistema interagiscono tra loro in modo non lineare e mutevole. Le relazioni non sono stabili.
		\item \textbf{Olismi ed Emergenza:} Il sistema \textbf{non è la semplice somma delle parti}. Dalle interazioni emergono comportamenti collettivi che non sono spiegabili analizzando i singoli componenti isolatamente (avanza sempre "qualcosa" nel passaggio dal micro al macro).
		\item \textbf{Prevedibilità limitata:} Un sistema complesso è solo parzialmente prevedibile. Piccole variazioni nelle condizioni iniziali possono produrre grandi cambiamenti nel risultato finale.
	\end{itemize}
	
	\paragraph{Esempio: Il valore sistemico di un Datacenter}
	Si consideri un impianto tecnologico, come un \textbf{datacenter}. Preso singolarmente sul mercato, esso ha un determinato prezzo di acquisto (valore di mercato). Tuttavia, nel momento in cui viene inserito in un'organizzazione specifica:
	\begin{itemize}
		\item Stabilisce relazioni con le altre risorse (capitale umano, software proprietari, flussi operativi).
		\item Il suo valore per l'azienda diventa diverso (spesso superiore) rispetto al semplice prezzo di mercato del bene isolato.
	\end{itemize}
	Questa differenza di valore è la prova che l'azienda è un \textbf{sistema complesso}: il valore generato deriva dalla qualità delle relazioni tra le parti e non solo dal valore dei singoli "pezzi" che la compongono. 
	
	\subsection{Cos'è un'impresa?}
	La nascita di un'impresa può essere ricondotta a due dinamiche principali:
	\begin{itemize}
		\item \textbf{Market Pull (Bisogno):} L'impresa nasce per rispondere a un'esigenza specifica già esistente nel mercato.
		\item \textbf{Technology/Product Push:} L'impresa nasce da un'innovazione o un prodotto sviluppato, con l'obiettivo di "creare" un nuovo mercato.
	\end{itemize}
	
	\subsubsection{Il finanziamento e le fonti di capitale}
	Per trasformare un'idea di business in realtà serve il \textbf{capitale}. Se l'imprenditore non ne possiede a sufficienza, deve rivolgersi all'esterno. Qui si distinguono due categorie fondamentali di finanziatori:
	
	\begin{enumerate}
		\item \textbf{Gli Azionisti (Soci):} Apportano il \textbf{Capitale di rischio} (o \textbf{Capitale proprio}). 
		Essendo "sulla stessa barca" dell'imprenditore, condividono il rischio d'impresa: se l'azienda fallisce, il loro investimento va perso.
		\item \textbf{I Creditori Finanziari (Banche):} Apportano il \textbf{Capitale di terzi} (o \textbf{Mezzi di terzi} / \textbf{Capitale di credito}). 
		Hanno i "piedi per terra": non sopportano il rischio operativo dell'impresa ma vantano un diritto legale alla restituzione della somma prestata.
	\end{enumerate}
	
	\subsubsection{Remunerazione del capitale}
	La distinzione tra queste due figure si riflette anche nel modo in cui vengono ricompensate:
	
	\begin{table}[h!]
		\centering
		\begin{tabular}{@{}lll@{}}
			\toprule
			\textbf{Tipo di Capitale} & \textbf{Chi lo apporta} & \textbf{Compenso (Remunerazione)} \\ \midrule
			\textbf{Proprio / Rischio} & Soci / Azionisti & \textbf{Utile} distribuito (Dividendi) \\
			\textbf{Di Terzi / Credito} & Banche / Finanziatori & \textbf{Interessi} (Costo del debito) \\ \bottomrule
		\end{tabular}
	\end{table}
	
	\begin{itemize}
		\item L'\textbf{Utile} è un compenso eventuale: i soci vengono pagati solo se l'azienda genera valore dopo aver coperto tutti i costi.
		\item Gli \textbf{Interessi} rappresentano invece il prezzo del denaro prestato: la banca deve riceverli a prescindere dal fatto che l'azienda sia in utile o in perdita.
	\end{itemize}
	
	\subsection{Il ciclo della gestione e la creazione di valore}
	L'impresa opera attraverso un ciclo continuo di attività. Inizialmente, vengono acquisite le risorse necessarie alla produzione:
	\begin{itemize}
		\item \textbf{Risorse materiali a fecondità ripetuta:} Immobili, impianti, stabilimenti (asset che durano nel tempo).
		\item \textbf{Risorse materiali a fecondità semplice:} Materie prime (che si esauriscono in un solo ciclo).
		\item \textbf{Risorse immateriali:} Software, brevetti, licenze.
		\item \textbf{Capitale Umano:} Assunzione di dipendenti e collaboratori.
	\end{itemize}
	
	Una volta trasformate le risorse e immessi i prodotti sul mercato, il denaro ricavato dalle vendite deve essere utilizzato secondo un ordine di priorità:
	\begin{enumerate}
		\item Copertura dei costi operativi e dei fornitori.
		\item \textbf{Pagamento degli interessi} ai creditori finanziari (banche).
		\item \textbf{Pagamento delle imposte} allo Stato.
		\item \textbf{Distribuzione dei dividendi} agli azionisti (solo se rimane un utile residuo).
	\end{enumerate}
	
	\subsubsection{La gestione come sistema dinamico}
	La \textbf{gestione d'impresa} può essere definita come un loop (circuito) continuo di operazioni di finanziamento, acquisto, produzione e vendita. 
	
	\begin{tcolorbox}[colback=green!5,colframe=green!60!black]
		\textbf{Condizione di sopravvivenza:} Il circuito continua a funzionare se, nel tempo, il \textbf{valore dell'output} (ciò che creiamo) è stabilmente \textbf{superiore al valore degli input} (risorse necessarie alla produzione). 
	\end{tcolorbox}
	
	\subsubsection{Il tempo e il bilancio}
	L'impresa è un fenomeno che si sviluppa nel \textbf{tempo}. Tuttavia, per esigenze pratiche e legali, abbiamo bisogno di misurarne l'andamento:
	\begin{itemize}
		\item \textbf{Il Bilancio:} È una "foto" scattata in un istante preciso (solitamente al 31/12) che cerca di riassumere cosa è accaduto in un determinato intervallo di tempo.
		\item \textbf{Risultato d'esercizio vs Valore reale:} Un risultato numerico negativo in bilancio (perdita) non coincide necessariamente con il fallimento imminente. Il valore di un'impresa si valuta sulla sua capacità di generare ricchezza nel lungo periodo, non solo su un singolo fermo immagine temporale.
	\end{itemize}
	
	\subsection{Il ruolo del Management: Arte e Scienza}
	La gestione d'impresa non può essere ridotta a una semplice sequenza di procedure o equazioni matematiche. Sebbene esistano tecniche rigorose, la natura di \textbf{sistema complesso} dell'azienda richiede un approccio che sia al contempo:
	\begin{itemize}
		\item \textbf{Scienza:} basata su modelli, analisi dei dati e processi standardizzati.
		\item \textbf{Arte:} legata all'intuito, alla capacità di adattamento e alla gestione delle relazioni umane in contesti imprevedibili.
	\end{itemize}
	
	\subsubsection{Efficienza vs Efficacia}
	Questi due concetti rappresentano la bussola del manager, ma misurano cose diverse:
	
	\begin{itemize}
		\item \textbf{Efficienza:} Misura la capacità di ottimizzare l'uso delle risorse. È il rapporto tra i risultati ottenuti (output) e le risorse impiegate (input).
		$$ \text{Efficienza} = \frac{\text{Risultato (monetario, qualitativo, ecc.)}}{\text{Risorse impiegate}} $$
		\textit{Essere efficienti significa "fare le cose nel modo giusto" (minimizzando gli sprechi).}
		
		\item \textbf{Efficacia:} Misura la capacità di raggiungere l'obiettivo prefissato. È il rapporto tra i risultati effettivamente ottenuti e quelli che si sarebbero potuti (o voluti) ottenere.
		$$ \text{Efficacia} = \frac{\text{Risultato Ottenuto}}{\text{Risultato Obiettivo}} $$
		\textit{Essere efficaci significa "fare le cose giuste" (raggiungere il traguardo).}
	\end{itemize}
	
	\subsubsection{I quattro pilastri del lavoro manageriale}
	Il processo di gestione si articola in quattro fasi fondamentali, spesso cicliche:
	
	\begin{enumerate}
		\item \textbf{Pianificazione:} Definizione degli obiettivi a lungo termine. Risponde alla domanda: "Cosa vogliamo che la nostra impresa diventi nel futuro?".
		\item \textbf{Organizzazione:} Definizione del \textit{come} raggiungere gli obiettivi. Riguarda il coordinamento e l'allocazione delle risorse (materiali, immateriali e persone) all'interno della struttura aziendale.
		\item \textbf{Leadership:} La componente umana del management. Consiste nel motivare, guidare e influenzare le persone affinché collaborino attivamente verso la visione comune.
		\item \textbf{Controllo:} Monitoraggio costante dei risultati. Consiste nel confrontare ciò che è stato realizzato con gli obiettivi pianificati, per correggere eventuali deviazioni.
	\end{enumerate}\subsection{Stakeholder e Shareholder}
	Nella gestione d'impresa è necessario distinguere tra due categorie di soggetti:
	
	\begin{itemize}
	\item \textbf{Stakeholder (Portatori di interesse):} Tutti i soggetti che hanno un interesse (economico, sociale o istituzionale) nell'andamento dell'impresa. 
	Tra questi troviamo: manager, lavoratori, fornitori, banche, Stato, clienti e perfino i concorrenti.
	\item \textbf{Shareholder (Azionisti):} Portatori di quote (\textit{shares}). Sono i soci proprietari, ovvero coloro che hanno conferito il capitale di rischio. Fanno parte della categoria degli stakeholder, ma con un legame di proprietà e rischio diretto.
	\end{itemize}
	
	\begin{quote}
	\textbf{Nota sulla forma giuridica:} Mentre la \textit{forma giuridica} definisce chi risponde legalmente delle obbligazioni dell'impresa, il concetto di \textit{stakeholder} definisce chi ha il diritto (o la necessità) di conoscere lo stato di salute dell'azienda.
	\end{quote}
	
	\subsection{Imprenditore vs. Manager}
	Sebbene in molte piccole imprese italiane queste figure coincidano, si tratta di ruoli con obiettivi e mentalità differenti.
	
	\begin{table}[h!]
		\centering
		\large
		\textbf{\textcolor{red}{IMPRENDITORE VS MANAGER: DUE RUOLI DIVERSI}}\\[0.5cm]
		\renewcommand{\arraystretch}{1.8}
		\begin{tabular}{l p{5.5cm} p{5.5cm}}
			\rowcolor{black}
			& \textbf{\textcolor{white}{IMPRENDITORE}} & \textbf{\textcolor{white}{MANAGER}} \\
			
			\textbf{\textcolor{red}{RUOLO}} & Crea l'impresa, assume il rischio, definisce la visione. & Gestisce l'impresa secondo gli obiettivi definiti. \\ \hline
			
			\textbf{\textcolor{red}{ORIZZONTE}} & Lungo termine (5-10+ anni). Pensa al futuro. & Medio-breve termine (1-3 anni). Pensa ai risultati. \\ \hline
			
			\textbf{\textcolor{red}{RISCHIO}} & Alto rischio. Scommette il proprio capitale. & Rischio moderato. Protegge il capitale altrui. \\ \hline
			
			\textbf{\textcolor{red}{MOTIVAZIONE}} & Visione, passione, creazione di valore. & Efficienza, profitti, stabilità. \\ \hline
		\end{tabular}
	\end{table}	
	
	\subsubsection{L'aspetto psicologico e temporale}
	L'imprenditore vive l'impresa come una propria estensione (una "creatura"). La sua preoccupazione è la continuità nel tempo. 
	Al contrario, il \textbf{manager} è un dipendente di alto livello: la sua priorità sono i risultati tangibili nel breve periodo (trimestrali, annuali), spesso perché la sua retribuzione e la sua carriera dipendono dal raggiungimento di target immediati. Questa differenza può generare tensioni se gli obiettivi di breve termine del manager sacrificano la visione di lungo termine dell'imprenditore.
	
	Short-Termismo: La tendenza a prendere delle decisioni che hanno un effetto nell'immediato, a discapito di decisione che vanno a impattare nel medio-lungo termine
	
	\section{Il caso Olivetti}
	\subsection{Caso di Studio: Il Gruppo Olivetti}
	La storia della Olivetti rappresenta un esempio emblematico della tensione tra visione imprenditoriale di lungo periodo e gestione manageriale focalizzata sul breve termine.
	
	\subsubsection{Dalla meccanica all'elettronica}
	\begin{itemize}
		\item \textbf{1908:} Camillo Olivetti fonda a Ivrea la prima fabbrica italiana di macchine da scrivere.
		\item \textbf{Anni '50:} Sotto la guida del figlio, \textbf{Adriano Olivetti}, l'azienda raggiunge la leadership mondiale (celebre il successo della \textit{Lettera 22}).
		\item \textbf{1954--1957:} Adriano ha una visione che anticipa i tempi: capisce che il futuro non è nella meccanica, ma nell'elettronica. Nasce la \textbf{Divisione Elettronica}.
		\item \textbf{1965:} Viene lanciato a New York il \textit{Programma 101}, considerato il primo vero "personal computer" al mondo.
	\end{itemize}
	
	\subsubsection{La crisi e il "Gruppo di Intervento"}
	Il cuore del caso esplode negli anni '60:
	\begin{itemize}
		\item \textbf{1960:} Muore improvvisamente Adriano Olivetti. L'azienda perde il suo leader visionario.
		\item \textbf{1960--1964:} Olivetti acquisisce l'americana \textbf{Underwood} per 40 milioni di dollari. L'operazione prosciuga la liquidità aziendale, rendendo i debiti insostenibili.
		\item \textbf{1964:} Per evitare il fallimento, interviene un consorzio di \textbf{stakeholder} (banche, Fiat, Pirelli), noto come "Gruppo di Intervento".
	\end{itemize}
	
	\subsubsection{Visione vs. Risultato immediato}
	La mentalità dei manager del Gruppo di Intervento era diametralmente opposta a quella di Adriano. Celebre rimane la loro dichiarazione sulla divisione elettronica:
	
	\begin{quote}
		\textit{"La divisione elettronica è un \textbf{tumore} che drena risorse vitali senza produrre utili immediati e va rimossa per salvare il core business meccanico."}
	\end{quote}
	
	Questo passaggio storico dimostra come:
	\begin{enumerate}
		\item L'imprenditore (Adriano) vedesse un \textbf{valore futuro} (l'elettronica) che ancora non esisteva nei bilanci.
		\item I manager (il Gruppo) si limitassero a guardare la \textbf{"foto" del bilancio} del momento, dove l'elettronica appariva solo come un costo, portandoli a voler amputare quella che sarebbe stata la più grande opportunità di crescita dell'azienda.
	\end{enumerate}
	
	\subsection{Il conflitto nel Board (1964): Lo scontro dei Trade-off}
	Nel 1964, il Consiglio di Amministrazione di Olivetti diventa il teatro di uno scontro ideologico e strategico tra due anime: i \textbf{Manager Finanziari} (esponenti del Gruppo di Intervento) e gli \textbf{Ingegneri Visionari} (gli eredi spirituali di Adriano).
	
	\begin{tcolorbox}[colback=red!5,colframe=red!75!black,title=Definizione di Trade-off]
		Il \textbf{Trade-off} (o relazione di scambio) è una situazione che implica una scelta tra due o più opzioni tra loro contrastanti: il miglioramento di un aspetto comporta necessariamente il peggioramento di un altro. In economia, rappresenta la rinuncia a un beneficio in cambio di un altro ritenuto più urgente o importante.
	\end{tcolorbox}
	
	\subsubsection{Le due posizioni a confronto}
	La battaglia si giocò su un trade-off fondamentale tra \textbf{sopravvivenza immediata} e \textbf{primato tecnologico futuro}:
	
	\begin{itemize}
		\item \textbf{I Manager Finanziari:}
		\begin{itemize}
			\item \textit{Logica:} Pragmatismo e avversione al rischio.
			\item \textit{Argomentazione:} "Dobbiamo vendere la divisione elettronica per risanare i conti. Il nostro dovere è proteggere i 50.000 posti di lavoro legati alla meccanica. L'elettronica è un settore troppo rischioso e incerto."
			\item \textit{Obiettivo:} Stabilità finanziaria di breve periodo.
		\end{itemize}
		
		\item \textbf{Gli Ingegneri Visionari:}
		\begin{itemize}
			\item \textit{Logica:} Innovazione e visione di lungo periodo.
			\item \textit{Argomentazione:} "Non possiamo vendere. Il nostro laboratorio ha quasi ultimato il \textbf{Programma 101}, il primo computer al mondo. Se rinunciamo all'elettronica, rinunciamo al nostro futuro e al primato mondiale."
			\item \textit{Obiettivo:} Creazione di valore e leadership tecnologica di lungo periodo.
		\end{itemize}
	\end{itemize}
	
	\subsubsection{Analisi del trade-off}
	
	Questo è un caso di trade-off classico.%
	\begin{enumerate}
		\item \textbf{Profitto immediato (e salvataggio):} garantito dal taglio dei costi (cessione della divisione), al prezzo però di perdere la capacità di innovare.
		\item \textbf{Profitto totale (e potenziale):} puntare sul rischio dell'elettronica, mettendo tuttavia a repentaglio la liquidità residua dell'azienda nel breve termine.
	\end{enumerate}
	
	Poi sappiamo tutti com'è finita: l'azienda ha venduto la parte elettronica a un'impresa statunitense, perdendo per sempre la sua leadership tecnologica.
	
	\paragraph{Trade-off}
	
	Un trade-off è una relazione di incompatibilità (totale o parziale) fra due obiettivi o variabili, tale per cui migliorare uno implica peggiorare l'altro a parità di risorse e tecnologia disponibili.\footnote{In economia, si parla di trade-off quando la crescita di una variabile risulta incompatibile con la crescita dell'altra e ne provoca una contrazione. } 
	Il trade-off rappresenta quindi un vincolo strutturale del sistema.%
	Un esempio semplice: una startup con 10 sviluppatori può allocare il tempo tra (1) aggiungere nuove funzionalità e (2) migliorare la qualità del codice.%
	Se privilegia (1), inevitabilmente non potrà migliorare (o potrebbe addirittura peggiorare) (2): c'è quindi un trade-off tra velocità e qualità, che può essere sia qualitativo (velocità–qualità) sia tecnico.
	
	Nel caso Olivetti, il trade-off fondamentale era tra \emph{valore} e \emph{profitto}.
	
	\paragraph{Costo-opportunità}
	
	Il costo-opportunità è il valore (espresso in termini economici) della migliore alternativa alla quale si rinuncia quando si prende una decisione.\footnote{Il costo-opportunità è il valore della migliore alternativa cui si rinuncia, cioè il costo della scelta non intrapresa. } 
	A differenza del trade-off, che descrive una tensione strutturale tra obiettivi incompatibili, il costo-opportunità è una misura quantitativa della rinuncia: se scelgo A invece di B, il costo-opportunità è il valore di B cui rinuncio scegliendo A.
	
	Esempio: una startup deve decidere tra (A) crescere rapidamente e (B) essere profittevole fin da subito.%
	Per crescere rapidamente occorrono investimenti che riducono i margini (ad esempio prezzi più bassi e forti spese commerciali).%
	\begin{itemize}
		\item Trade-off 1: se aumento la crescita, devo accettare margini più bassi.
		\item Trade-off 2: se aumento i margini (prezzi più alti, meno investimenti), venderò meno e crescerò più lentamente.
	\end{itemize}
	Il costo-opportunità della scelta di crescita aggressiva è il profitto cui rinuncio rispetto a una strategia di crescita più lenta ma più redditizia nel breve termine.
	
	\section{Caso Ferrari}
	
	Il caso Ferrari ha un esito molto diverso rispetto a Olivetti.%
	Ferrari decide di realizzare un'auto elettrica, pur essendo un'azienda che potrebbe vendere molte più auto l'anno rispetto alle circa 13\,000 effettivamente prodotte.%
	Come spiegava Enzo Ferrari: «Venderò sempre un'auto in meno di quante il mercato ne richiede.»%
	Questo garantisce che il desiderio per il brand rimanga sempre parzialmente insoddisfatto: la lista d'attesa diventa un vero e proprio rito di passaggio, rafforzando lo status del cliente e consolidando il valore simbolico del marchio.%
	La scarsità di auto contribuisce inoltre a proteggere il valore dell'usato, rendendo una Ferrari percepita come un investimento e non solo come un bene di consumo.
	
	A questo punto si arriva al \emph{dilemma 2026}, che procede in parallelo con le decisioni europee sui motori termici (divieto di vendita di vetture nuove con motore a combustione dal 2035).\footnote{Il phase-out dei motori a combustione interna per le auto nuove entro il 2035 è parte del pacchetto europeo per la transizione ecologica. } 
	Ha senso una Ferrari senza il suo iconico rombo?%
	Questo apre una serie di sfide tecniche e di brand:
	\begin{enumerate}
		\item progettare un suono sintetico coerente con l'identità Ferrari;
		\item riprodurre il feedback tattile (vibrazioni e risposta fisica del motore);
		\item sviluppare un'interfaccia uomo–macchina che trasmetta l'eredità del mondo Ferrari con un'esperienza immersiva;
		\item gestire il peso delle batterie garantendo l'agilità tipica del marchio.
	\end{enumerate}
	
	Nel 2024 emergono tre vincoli principali:
	\begin{enumerate}
		\item \textbf{Vincolo normativo:} il quadro regolatorio europeo vieta la vendita di nuove auto con motore a combustione interna dal 2035.
		\item \textbf{Vincolo finanziario:} la transizione all'elettrico richiede investimenti per miliardi; questo deve essere conciliato con gli interessi degli azionisti, che desiderano ritorni elevati e rapidi. La funzione obiettivo dell'azionista può condizionare le scelte manageriali, orientando il management verso decisioni che massimizzano il breve periodo a scapito del medio–lungo termine.
		\item \textbf{Vincolo di brand:} Ferrari vende esclusività. Se aumenta troppo i volumi produttivi rischia di perdere l'esclusività; se li mantiene troppo bassi, potrebbe non coprire i costi fissi e gli investimenti in R\&S.
	\end{enumerate}
	
	Nel 2024 il Consiglio di amministrazione si pone quindi il problema di:
	\begin{itemize}
		\item investire nell'elettrico,
		\item mantenere l'identità Ferrari che conosciamo.
	\end{itemize}
	
	Si confrontano due proposte estreme:
	\begin{description}
		\item[Proposta A -- Esclusività:] volumi molto ridotti (ad esempio, 2\,000 unità all'anno) con un prezzo di circa 500\,000 euro per auto.
		\item[Proposta B -- Leadership di mercato:] volumi più elevati (ad esempio, 10\,000 unità all'anno) con un prezzo di circa 200\,000 euro, puntando sulla quota di mercato.
	\end{description}
	
	C'è chi sostiene: «meno auto a un prezzo più elevato» e chi, al contrario, «più auto a un prezzo più basso».%
	La strategia effettivamente adottata è una via di mezzo: circa 5\,000 unità all'anno a un prezzo di 400\,000 euro.%
	Gli effetti di questa scelta si vedranno nel medio periodo, ma la logica dichiarata è: «produciamo abbastanza volume per finanziare le spese di R\&S sull'elettrico, ma non così tanto da rendere l'auto alla portata di tutti».%
	La preoccupazione principale resta quindi la protezione del brand.
	
	Anche in questo caso, come in quello Olivetti, emergono diversi trade-off:
	\begin{enumerate}
		\item \textbf{Volumi (quantità) vs margine:} più aumento i volumi, più devo contenere i prezzi; se alzo i prezzi per difendere l'esclusività, devo accettare volumi più bassi.
		\item \textbf{Motore termico vs transizione all'elettrico:} come mantenere l'emozione Ferrari nel passaggio a una tecnologia radicalmente diversa.
		\item \textbf{Performance vs comfort:} bilanciare prestazioni estreme con comfort e fruibilità quotidiana.
		\item \textbf{Imprenditore vs manager:} tensione ricorrente tra visione di lungo periodo (tipica dell'imprenditore–fondatore) e vincoli di breve periodo legati a governance e mercato dei capitali.
	\end{enumerate}
	
	\section{Elementi di disturbo del trade-off}
	
	I trade-off sono vincoli strutturali del sistema.%
	Esistono però anche elementi di \emph{disturbo}, ossia situazioni organizzative o di fatto che trasformano un obiettivo comune in una serie di «o questo o quello» operativi.
	
	Esempio: all'inizio dell'anno viene definito il \emph{budget}, ossia il piano degli obiettivi per i successivi 10--12 mesi.%
	Il direttore generale convoca i responsabili di funzione e comunica l'obiettivo annuale dell'impresa: «Aumentare la nostra quota di mercato».%
	Ciascun direttore interpreta questo obiettivo attraverso le «lenti» della propria funzione:
	\begin{itemize}
		\item il direttore di produzione lo traduce in «produrre di più»;
		\item il direttore marketing lo declina in «migliorare l'immagine del brand»;
		\item il direttore vendite lo interpreta come «vendere di più».
	\end{itemize}
	Tutti e tre rispondono al macro–obiettivo, ma a livello operativo nascono possibili conflitti e vincoli.
	
	Ogni direttore, a sua volta, deve tradurre l'obiettivo della funzione in azioni concrete per il proprio team.%
	Per esempio:
	\begin{itemize}
		\item produzione può proporre di investire in un nuovo stabilimento o di ricorrere alla subfornitura;
		\item marketing può pianificare una campagna pubblicitaria intensiva per rafforzare il posizionamento di qualità;
		\item vendite può proporre politiche di sconto aggressive per aumentare i volumi.
	\end{itemize}
	Tra queste azioni possono emergere trade-off: ad esempio, la scelta di vendite di ridurre i prezzi può entrare in tensione con il posizionamento di alta qualità costruito dal marketing (che richiede investimenti pubblicitari da recuperare attraverso margini più elevati).%
	Anche quando l'obiettivo strategico è chiaro, i sotto–obiettivi funzionali possono creare interdipendenze e vincoli che generano ulteriori trade-off operativi.
	
	\section{Caso Nokia}
	
	Fino al 2007 Nokia era al vertice della telefonia mobile globale.%
	Aveva lanciato il primo smartphone nel 1996 e il primo telefono con fotocamera nel 2001.\footnote{Nokia fu tra i primi produttori a introdurre smartphone e fotocamere sui telefoni mobili.} 
	Nel 2007 deteneva oltre il 40\% della quota del mercato globale e ricavi superiori ai 50 miliardi di euro, con una crescita dell'utile operativo intorno al 45\%.%
	La quota di mercato del sistema operativo proprietario Symbian era circa il 60\% (sistema chiuso e altamente ottimizzato per l'hardware Nokia).\footnote{Symbian OS era il sistema proprietario di Nokia, dominante nel mercato smartphone prima del 2007. } 
	
	A gennaio 2007 esce il primo iPhone.%
	Quando l'iPhone arriva sul mercato, i principali competitor sono BlackBerry (segmento business) e Nokia (consolidata leader di mercato).%
	Fra il 2007 e il 2011 Nokia compie tre scelte che si riveleranno critiche:
	\begin{enumerate}
		\item decide di mantenersi fedele a Symbian;
		\item mostra resistenza al touchscreen, continuando a puntare sulla tastiera fisica mentre il mercato, in poco più di un anno, converge verso il touchscreen; il primo Nokia completamente touchscreen arriva solo a fine 2008;
		\item nel 2011 il CEO sceglie Windows Phone come piattaforma principale per gli smartphone Nokia.\footnote{L'11 febbraio 2011 Nokia annuncia una partnership strategica con Microsoft, adottando Windows Phone come sistema operativo principale sui propri smartphone. }
	\end{enumerate}
	La scommessa su Windows Phone si rivela fallimentare: Nokia perde ulteriori quote di mercato e la leadership nel segmento smartphone.\footnote{Windows Phone non riesce a colmare il divario con Android e iOS, e la quota di mercato di Nokia continua a ridursi. }
	
	\subsection*{Domande}
	
	\paragraph{1) Trade-off affrontati da Nokia (2007–2011)}
	
	Principali trade-off e relative scelte,rinunce:
	\begin{enumerate}
		\item \textbf{Hardware: tastiera fisica vs touchscreen}
		\begin{itemize}
			\item \emph{Scelta}: mantenere la tastiera fisica come elemento distintivo per troppo tempo.
			\item \emph{Rinuncia}: posizionarsi subito sullo standard emergente del touchscreen, rallentando l'adattamento al nuovo paradigma d'uso introdotto da iPhone e Android.
		\end{itemize}
		\item \textbf{Sistema operativo: Symbian (chiuso) vs Android (aperto)}
		\begin{itemize}
			\item \emph{Scelta}: continuare a investire su Symbian, sistema proprietario chiuso e ottimizzato per l'hardware Nokia.
			\item \emph{Rinuncia}: entrare rapidamente nell'ecosistema Android, beneficiando della crescita della piattaforma aperta e del supporto degli sviluppatori di terze parti.\footnote{Il concetto di trade-off tra sistemi chiusi e aperti è centrale nelle strategie di piattaforma.}
		\end{itemize}
		\item \textbf{Piattaforma strategica: Windows Phone vs Android}
		\begin{itemize}
			\item \emph{Scelta}: partnership con Microsoft e adozione di Windows Phone come piattaforma principale.
			\item \emph{Rinuncia}: allinearsi all'ecosistema Android, che nel frattempo diventava lo standard dominante (come dimostrato dalla crescita di player come Samsung).%
		\end{itemize}
	\end{enumerate}
	
	\paragraph{2) Costo-opportunità della mancata adozione di Android}
	
	Il costo-opportunità della decisione di non adottare Android nel 2008–2009 può essere rappresentato:
	\begin{itemize}
		\item dalla perdita di quota di mercato potenziale che Nokia avrebbe potuto ottenere entrando in Android nelle sue fasi iniziali;
		\item dalla rinuncia al mantenimento della leadership nel settore smartphone, ruolo che è stato invece assunto da produttori come Samsung, passata da una quota modesta nel 2009 a una quota molto elevata nel 2012 grazie proprio ad Android.\footnote{La crescita di Samsung nel mercato smartphone è strettamente legata all'adozione tempestiva di Android. }
	\end{itemize}
	In sintesi, scegliendo Android, Nokia avrebbe potuto beneficiare di una traiettoria di crescita simile a quella dei principali concorrenti che hanno abbracciato la piattaforma, preservando parte della sua posizione dominante.
	
	\paragraph{3) Costo-opportunità degli investimenti in Symbian}
	
	Nel 2007 Nokia aveva ricavi pari a circa 52 miliardi di euro.%
	Ogni euro investito in Symbian è, di fatto, un euro in meno investito nello sviluppo o nell'adozione di altri sistemi operativi (come Android o nuove piattaforme interne).%
	Il costo-opportunità qui è in larga parte \emph{implicito}: Nokia non prende la decisione esplicita di «spendere X per Android», ma di «continuare a investire su Symbian», ignorando il valore delle alternative.
	
	Questa distinzione è importante per almeno due motivi:
	\begin{enumerate}
		\item \textbf{Focalizzazione sui costi espliciti:} i manager tendono a concentrarsi su ciò che spendono (costi espliciti), trascurando ciò che non stanno spendendo (costi impliciti), come gli investimenti mancati su Android o su piattaforme alternative.
		\item \textbf{Vittima del proprio successo:} nel 2007 Nokia è leader di mercato, con profitti elevati e costi sotto controllo. Il costo del «restare su Symbian» non è percepito come un esborso immediato (costo esplicito), ma come una mancata opportunità futura; proprio per questo viene sottovalutato, contribuendo a ritardare la reazione strategica al cambiamento tecnologico.
	\end{enumerate}
	Se il costo del restare su Symbian fosse stato reso più esplicito (ad esempio tramite scenari di confronto strutturati con Android), la percezione del rischio e le decisioni di investimento sarebbero potute essere diverse.
	
	
	\subsection{Bias cognitivi}
	
	L'economia classica modella il decisore come un soggetto perfettamente razionale (\emph{homo oeconomicus}) che:%
	\begin{enumerate}
		\item conosce tutte le alternative disponibili sul mercato;
		\item è in grado di valutare correttamente tutti i costi-opportunità;%
		\item sceglie sempre l'opzione che massimizza la propria utilità;
		\item non è influenzato dal modo in cui le informazioni vengono presentate (\emph{framing}).%
	\end{enumerate}
	La realtà è diversa: ciascuno di noi è soggetto a \emph{bias cognitivi}, ossia pattern di deviazione sistematica rispetto al comportamento razionale atteso.%
	Questi bias distorcono la percezione dei trade-off e portano a sottovalutare o ignorare i costi-opportunità.%
	Nota: un bias non è un semplice errore casuale, ma una tendenza stabile e prevedibile.
	
	\subsubsection{Razionalità limitata (Herbert Simon)}
	
	Herbert A.~Simon (Premio Nobel per l'economia, 1978) introduce il concetto di \emph{razionalità limitata} (\emph{bounded rationality}).%
	Il processo decisionale è vincolato da tre fattori fondamentali:\footnote{Si veda la discussione sulla razionalità limitata nelle slide. }%
	\begin{enumerate}
		\item \textbf{Capacità cognitiva:} la mente umana non può elaborare simultaneamente tutte le informazioni rilevanti.
		\item \textbf{Informazioni disponibili:} le informazioni a disposizione sono sempre incomplete o imperfette, anche dopo ricerche approfondite.
		\item \textbf{Tempo:} il tempo a disposizione è limitato e ha un costo; non è possibile analizzare exhaustivamente tutte le opzioni.
	\end{enumerate}
	
	In questo contesto, gli individui non \emph{ottimizzano} ma \emph{satisfano} (\emph{satisficing}): cercano la prima soluzione \emph{abbastanza buona} che soddisfa soglie minime accettabili, interrompendo la ricerca di alternative potenzialmente migliori.\footnote{L'esempio in slide sul team che seleziona un fornitore mostra come si valuti solo un sottoinsieme di opzioni e ci si fermi al primo candidato adeguato. }%
	
	\subsubsection{Opportunity cost neglect (Frederick et al., 2009)}
	
	Frederick et al.~(2009) mostrano che le persone ignorano sistematicamente i costi-opportunità, introducendo il concetto di \emph{opportunity cost neglect}.\footnote{L'opportunity cost neglect è uno dei risultati più robusti dell'economia comportamentale: le alternative non scelte non vengono considerate spontaneamente.}%
	In pratica, gli individui:
	\begin{itemize}
		\item non considerano spontaneamente le alternative a cui rinunciano;
		\item non le vedono, non le percepiscono e quindi non le registrano mentalmente.
	\end{itemize}
	
	\paragraph{Esperimento del DVD Sony}
	
	\begin{itemize}
		\item \textbf{Condizione controllo:} «Compreresti il DVD Sony a 99\$?» (Opzioni: \emph{Sì} / \emph{No}).
		\item \textbf{Condizione reminder:} «Compreresti il DVD Sony a 99\$?» (Opzioni: \emph{Sì} / \emph{No, tenendo i 99\$ per altri acquisti}).%
	\end{itemize}
	La semplice aggiunta della frase «tenendo i 99\$ per altri acquisti» (informazione logicamente ridondante) riduce in modo significativo la percentuale di acquisto: rendere esplicita l'alternativa (\emph{non comprare e usare i soldi diversamente}) attiva la considerazione del costo-opportunità.\footnote{Plantinga et al.~(2017) replicano l'effetto su larga scala, mostrando che è presente sia in soggetti ad alto sia a basso reddito. }%
	
	\subsubsection{Contabilità mentale (Richard Thaler, 1985)}
	
	Richard Thaler (1985) introduce il concetto di \emph{contabilità mentale} (\emph{mental accounting}) per descrivere come le persone organizzano e gestiscono le proprie finanze.\footnote{La contabilità mentale porta a trattare somme di denaro identiche in modo diverso a seconda della loro «etichetta» mentale.}%
	Invece di considerare il denaro come perfettamente fungibile, gli individui creano \emph{conti mentali} separati (stipendio, bonus, risparmi, ecc.) e si comportano diversamente a seconda del conto.
	
	Una conseguenza importante è che i costi-opportunità vengono spesso sottovalutati rispetto ai costi espliciti:
	\begin{itemize}
		\item i costi espliciti (esborso di cassa visibile) sono immediatamente percepiti;
		\item i costi impliciti (alternative a cui si rinuncia) rimangono invisibili nella contabilità mentale.%
	\end{itemize}
	Questo è collegato al caso Nokia: l'investimento su Symbian era un costo esplicito, mentre il mancato investimento su Android era un costo-opportunità implicito, facilmente trascurato.
	
	\subsubsection{Avversione alle perdite e teoria del prospetto}
	
	Kahneman e Tversky (1979) sviluppano la \emph{Prospect Theory}, che rivoluziona la comprensione delle decisioni in condizioni di rischio.\footnote{La Prospect Theory mostra che la funzione di valore è concava nei guadagni e convessa nelle perdite, con un peso maggiore attribuito alle perdite. }%
	Uno dei risultati chiave è l'\emph{avversione alle perdite} (\emph{loss aversion}):
	\begin{itemize}
		\item le perdite hanno un impatto emotivo circa due volte superiore ai guadagni equivalenti;
		\item rinunciare a un attributo percepito come perdita pesa molto di più che acquisire un attributo percepito come guadagno.
	\end{itemize}
	
	Questo influenza direttamente i trade-off: nella valutazione di un investimento (ad esempio l'adozione di un nuovo sistema ERP) la «perdita» del vecchio sistema noto viene sovrastimata rispetto ai benefici del nuovo sistema più efficiente, portando spesso a ritardi nell'adozione di tecnologie superiori. %
	
	\subsubsection{Framing effect}
	
	Il \emph{framing effect} descrive come la stessa informazione, presentata in modo diverso, possa portare a decisioni differenti.\footnote{Tversky e Kahneman (1981) lo illustrano con il celebre esperimento dell'\"Asian Disease Problem\". }%
	L'informazione oggettiva non cambia, ma cambia la \emph{cornice} linguistica:
	
	\begin{itemize}
		\item \textbf{Frame positivo (guadagno):} «Un programma salverà 200 persone su 600.»
		\item \textbf{Frame negativo (perdita):} «Un programma causerà la morte di 400 persone su 600.»
	\end{itemize}
	I due enunciati sono logicamente equivalenti, ma nel frame di guadagno la maggioranza sceglie l'opzione sicura, mentre nel frame di perdita la maggioranza preferisce l'opzione rischiosa.%
	
	\paragraph{Esempio gestionale}
	
	Due presentazioni dello stesso progetto di automazione:
	\begin{itemize}
		\item «L'investimento permetterà di risparmiare 300.000 euro/anno in costi di manodopera.» (frame guadagno)
		\item «Se non investiamo, perderemo 300.000 euro/anno in inefficienze produttive.» (frame perdita)
	\end{itemize}
	La formulazione in termini di \emph{perdita evitata} genera una maggiore propensione all'investimento, a causa dell'avversione alle perdite. Chi presenta un progetto deve esserne consapevole.
	
	\subsubsection{Troppe alternative e paralisi decisionale}
	
	Un ulteriore effetto rilevante è la \emph{paralisi decisionale} dovuta all'eccesso di alternative (\emph{choice overload}).\footnote{Nelle slide viene suggerito di strutturare le opzioni gerarchicamente e di valutare 3–5 alternative per fase per mitigare il problema. }%
	Quando il numero di opzioni è troppo elevato:
	\begin{itemize}
		\item aumenta il carico cognitivo necessario per confrontare i trade-off tra alternative;
		\item cresce il timore di rimpiangere la scelta (\emph{regret}), con conseguente procrastinazione o non-decisione;
		\item il decisore tende a rifugiarsi nello \emph{status quo} o in euristiche semplici (prima opzione accettabile, scelta default, ecc.).
	\end{itemize}
	
	In contesti aziendali (progettazione, portafoglio progetti, selezione fornitori) questo implica la necessità di:
	\begin{itemize}
		\item ridurre e strutturare l'insieme delle alternative;
		\item lavorare per fasi, scremando progressivamente le opzioni;
		\item rendere espliciti i trade-off chiave invece di presentare liste lunghe e indifferenziate di scenari.%
	\end{itemize}
	\section{Le forme giuridiche delle imprese}
	
	\subsection{Introduzione: la scelta strategica della forma giuridica}
	
	Partendo da un'idea brillante e 1.000 euro in tasca, quale è la «strada legale» migliore per lanciare una startup tecnologica in Italia?%
	La scelta della forma giuridica non è una questione meramente burocratica, ma una decisione strategica cruciale che:
	\begin{itemize}
		\item determina chi paga i debiti aziendali;
		\item definisce la governance e i meccanismi decisionali;
		\item influenza la capacità di attrarre investimenti;
		\item accelera o blocca la crescita dell'impresa.
	\end{itemize}
	
	\subsubsection{Il caso CodiceSemplice: da freelance a imprenditore}
	
	Marco avvia la sua attività di sviluppo software come impresa individuale (Anno 1): costi minimi, test del mercato, piena autonomia decisionale, ma tutto il patrimonio personale è a rischio.%
	Al terzo anno, il successo cresce ma emergono problemi: non può assumere dipendenti qualificati, non può attrarre investitori (impossibile emettere quote), la casa è a rischio per i debiti aziendali.%
	Al quarto anno, Marco trasforma l'impresa in SRL: separa il patrimonio personale da quello aziendale, assume dipendenti e fa entrare un socio investitore.
	
	\subsection{Concetti fondamentali}
	
	\paragraph{Impresa}
	
	Attività economica organizzata e professionale volta a produrre o scambiare beni e servizi.
	
	\paragraph{Ditta}
	
	Nome commerciale dell'imprenditore individuale.%
	Deve contenere almeno il cognome del titolare.
	
	\paragraph{Società}
	
	Contratto tra due o più persone che conferiscono beni o servizi per esercitare insieme un'attività economica, allo scopo di dividerne gli utili.
	
	\paragraph{Distinzione chiave}
	
	L'impresa individuale non ha personalità giuridica separata dal titolare: l'imprenditore e l'impresa coincidono.%
	Le società di capitali, invece, godono di personalità giuridica autonoma: la società è un soggetto distinto dai soci.
	
	\subsection{L'impresa individuale}
	
	\subsubsection{Caratteristiche}
	
	\begin{itemize}
		\item \textbf{Unico titolare:} piena autonomia decisionale.
		\item \textbf{Nessun capitale minimo richiesto:} accessibile a tutti.
		\item \textbf{Costituzione rapida:} basta la Comunicazione Unica al Registro delle Imprese.
		\item \textbf{Ideale per:} freelance, test di mercato, attività di piccole dimensioni.
	\end{itemize}
	
	\subsubsection{Il grande rischio: responsabilità illimitata}
	
	L'imprenditore risponde dei debiti aziendali con \emph{tutto} il proprio patrimonio personale: casa, automobile, risparmi, conto corrente.%
	Ogni investimento mette a repentaglio il patrimonio privato.
	
	\subsubsection{I sei limiti alla crescita}
	
	\begin{enumerate}
		\item \textbf{Responsabilità illimitata:} il titolare risponde con tutti i beni personali; ogni investimento mette a rischio il patrimonio privato.
		\item \textbf{Accesso limitato al capitale:} impossibile emettere quote o azioni; niente venture capital, investitori o quotazione in borsa; crescita legata solo all'autofinanziamento.
		\item \textbf{Nessuna personalità giuridica:} l'impresa coincide con il titolare; malattia o morte bloccano l'attività; la continuità aziendale è strutturalmente fragile.
		\item \textbf{Difficoltà a coinvolgere talenti:} impossibile offrire partecipazioni a dipendenti o manager; difficile attrarre figure chiave indispensabili per scalare il business.
		\item \textbf{Credibilità limitata:} grandi clienti e partner internazionali preferiscono interlocutori con struttura societaria; l'impresa individuale è percepita come realtà minima.
		\item \textbf{Carico fiscale penalizzante:} i redditi confluiscono nell'IRPEF personale (aliquota fino al 43\%); le SRL scontano l'IRES al 24\%, con pianificazione fiscale più efficiente.
	\end{enumerate}
	
	Soluzione naturale: trasformare l'impresa in SRL per limitare la responsabilità, attrarre investitori e garantire continuità.
	
	\subsection{Società di persone}
	
	Le società di persone si fondano sull'elemento personale (fiducia reciproca tra soci) e sono caratterizzate da \emph{autonomia patrimoniale imperfetta}: il patrimonio dei soci non è completamente separato da quello della società.
	
	\subsubsection{Caratteristiche comuni}
	
	\begin{itemize}
		\item \textbf{Responsabilità illimitata e solidale:} i soci rispondono con tutto il proprio patrimonio per le obbligazioni sociali; i creditori possono rivalersi sui beni personali dei soci (dopo aver tentato di soddisfarsi sul patrimonio sociale).
		\item \textbf{Non hanno personalità giuridica autonoma.}
		\item \textbf{Esiste il socio d'opera:} possibilità di conferire lavoro, competenze e know-how (oltre a denaro e beni).
		\item \textbf{Non esistono organi sociali formali:} i soci sono anche amministratori.
	\end{itemize}
	
	\subsubsection{Società in Nome Collettivo (SNC)}
	
	\begin{itemize}
		\item Tutti i soci hanno responsabilità illimitata e solidale (articolo 2291 Codice Civile).
		\item Tutti i soci hanno, in linea di principio, potere di amministrazione.
		\item Nessun capitale minimo richiesto.
		\item Ideale per: piccole attività tra soci con forte legame fiduciario.
	\end{itemize}
	
	\paragraph{Responsabilità solidale}
	
	Anche un socio con una minima quota può essere obbligato a pagare l'intero debito contratto dalla società o da un altro socio.%
	Il comportamento scorretto di un socio (illeciti fiscali, abusi contrattuali, errori gestionali) ricade su tutti.
	
	\subsubsection{Società in Accomandita Semplice (SAS)}
	
	La SAS prevede due categorie di soci con ruoli e responsabilità distinte:
	
	\begin{table}[h]
		\centering
		\begin{tabular}{|l|l|}
			\hline
			\textbf{Accomandatari} & \textbf{Accomandanti} \\
			\hline
			Gestiscono l'azienda & Solo finanziatori \\
			Responsabilità illimitata e solidale & Responsabilità limitata al capitale conferito \\
			Amministrano la società & Esclusi dall'amministrazione \\
			\hline
		\end{tabular}
	\end{table}
	
	Ideale per situazioni in cui alcuni soci vogliono investire senza gestire e senza rischiare tutto il patrimonio personale.
	
	\subsection{Società di capitali}
	
	Le società di capitali si fondano sull'elemento patrimoniale (capitale sociale) e godono di \emph{autonomia patrimoniale perfetta}: il patrimonio della società è completamente separato da quello dei soci.%
	La società risponde delle obbligazioni sociali solo con il proprio patrimonio.
	
	\subsubsection{Caratteristiche comuni}
	
	\begin{itemize}
		\item \textbf{Responsabilità limitata:} i soci rischiano solo il capitale investito; i beni personali (casa, auto, risparmi) restano al sicuro dai creditori aziendali.
		\item \textbf{Hanno personalità giuridica autonoma:} la società è un soggetto di diritto distinto dai soci.
		\item \textbf{Non esiste il socio d'opera:} i conferimenti devono essere in denaro o beni valutabili economicamente (non è possibile conferire solo lavoro).
		\item \textbf{Organizzazione corporativa con organi sociali:}
		\begin{itemize}
			\item Assemblea dei soci (organo deliberativo)
			\item Amministratori (singolo o Consiglio di Amministrazione; può esserci un amministratore delegato)
			\item Collegio Sindacale (organo di controllo, obbligatorio in alcuni casi)
		\end{itemize}
	\end{itemize}
	
	\subsubsection{Società a Responsabilità Limitata (SRL)}
	
	\begin{itemize}
		\item \textbf{Capitale minimo:} 10.000 euro (Eccezione: 1 euro per SRLS semplificata).%
		Se il capitale è inferiore a 10.000 euro, i conferimenti devono essere versati interamente in denaro al momento della sottoscrizione.%
		Se il capitale è pari o superiore a 10.000 euro, è sufficiente versare il 25\% in denaro e il 100\% dei beni conferiti.
		\item \textbf{Partecipazioni:} quote (non azioni, non hanno un mercato e non circolano in borsa).
		\item \textbf{Governance:} Semplificata, flessibile, adattabile alle esigenze dei soci.
		\item \textbf{Collegio Sindacale:} obbligatorio solo se la SRL supera per due esercizi consecutivi due delle seguenti soglie:
		\begin{itemize}
			\item almeno 4.400.000 euro di attivo di stato patrimoniale;
			\item almeno 8.800.000 euro di ricavi da vendite e prestazioni;
			\item almeno 50 dipendenti occupati in media durante l'esercizio.
		\end{itemize}
		Oppure se la SRL è obbligata alla redazione del bilancio consolidato o controlla società soggette a revisione legale.
		\item \textbf{Ideale per:} PMI, startup, imprese familiari che vogliono limitare il rischio e mantenere flessibilità gestionale.
	\end{itemize}
	
	\subsubsection{Società per Azioni (SPA)}
	
	\begin{itemize}
		\item \textbf{Capitale minimo:} 50.000 euro (ridotto da 120.000 euro con il D.L. 91/2014).%
		Alla sottoscrizione dell'atto costitutivo deve essere versato in denaro almeno il 25\% dei conferimenti.
		\item \textbf{Partecipazioni:} azioni (liberamente trasferibili, salvo clausole statutarie).
		\item \textbf{Governance:} strutturata e formale; presenza obbligatoria di organi sociali (assemblea, consiglio di amministrazione, collegio sindacale).
		\item \textbf{Ideale per:} grandi imprese, società quotate in borsa (IPO), imprese che necessitano di raccogliere ingenti capitali dal mercato pubblico.
	\end{itemize}
	
	\subsubsection{Società in Accomandita per Azioni (SAPA)}
	
	Modello ibrido: struttura da SPA, anima da società di persone.
	
	\begin{itemize}
		\item \textbf{Soci accomandatari:} sono amministratori di diritto; responsabilità illimitata e solidale. Sono soci del tutto identici a quelli della società di persone. Loro responsabilità va ben oltre alla quota conferita (l'azione), si estende all'intero loro patrimonio.
		\item \textbf{Soci accomandanti:} esclusi dall'amministrazione; responsabilità limitata al capitale. Possiedono azioni!
	\end{itemize}
	
	Utilizzo in Italia: modello raro, utilizzato principalmente dalle grandi dinastie imprenditoriali per mantenere il controllo familiare stabile e indissolubile (es. Giovanni Agnelli B.V. per la famiglia Agnelli).
	
	\subsection{Startup innovative}
	
	La startup innovativa non è una forma giuridica, ma uno \emph{status} (D.L. 179/2012) applicabile a società di capitali (SRL o SPA) che rispettino determinati requisiti. 
	
	\subsubsection{Requisiti principali}
	I requisiti sono 3 più c'è un criterio.
	\begin{itemize}
		\item Costituita da meno di 5 anni.
		\item Sede legale in Italia (o in altro Paese UE purché con sede produttiva o filiale in Italia).
		\item Fatturato annuo inferiore a 5 milioni di euro.
		\item Non distribuisce utili.
		\item Oggetto sociale: sviluppo, produzione e commercializzazione di prodotti o servizi innovativi ad alto valore tecnologico.
		\item Deve soddisfare almeno uno dei seguenti criteri:
		\begin{itemize}
			\item spese in R\&S (investimento in Ricerca e Sviluppo) pari ad almeno il 15\% del maggiore tra costo e valore della produzione;
			\item almeno 1/3 del personale con dottorato di ricerca o 2/3 con laurea magistrale;
			\item titolarità di brevetti, software registrati o licenze su proprietà industriale.
		\end{itemize}
	\end{itemize}
	
	\subsubsection{Vantaggi concreti}
	
	\begin{itemize}
		\item \textbf{Zero burocrazia:} esenzione da diritti camerali e imposte di bollo.
		\item \textbf{Accesso al credito:} accesso gratuito e semplificato al Fondo di Garanzia per le PMI.
		\item \textbf{Incentivi fiscali:} detrazioni fiscali del 30--50\% per chi investe nel capitale sociale (regime «de minimis» con detrazione fino al 65\%).
		\item \textbf{Work for equity:} possibilità di remunerare collaboratori, consulenti e dipendenti con quote societarie, con regime fiscale agevolato.
		\item \textbf{Deroghe alla disciplina societaria ordinaria:} maggiore flessibilità nelle clausole statutarie.
		\item \textbf{Raccolta di capitali tramite equity crowdfunding.}
		\item \textbf{Servizi di internazionalizzazione} (ICE -- Agenzia per la promozione all'estero).
	\end{itemize}
	
	\subsubsection{Esempi di startup innovative italiane di successo}
	
	\begin{itemize}
		\item \textbf{Exein} (Cybersecurity IoT): fondata nel 2018, protegge dispositivi connessi; ha raccolto 100 milioni di euro per l'espansione globale.
		\item \textbf{Nanophoria} (Biotech): sviluppa trattamenti non invasivi per il cuore; ha raccolto 83,5 milioni di euro in un round di Serie A.
		\item \textbf{WSense} (Deep Tech): spin-off de La Sapienza, ha sviluppato il «wifi sottomarino» per il monitoraggio ambientale.
	\end{itemize}
	
	\begin{comment}%% NUOVO DEL 25 MARZO
	La lezione: le competenze ingegneristiche sono il cuore, ma la forma giuridica corretta è il motore che ha permesso a queste aziende di attrarre capitali e crescere.
	
	Cosa sono Spin-off Split-off e Split-up?
	Facciamo una premessa, quando parliamo di Spin-Off, Split-Up parliamo di natura straordinarie (non quotidiane), queste operazioni le possiamo collocare in una matrice 2x2, dove nelle righe si dice chi decide (Chi decide l'operazione? Può essere decisa su base volontaria o cessazione coatta "imposta dalla legge, forzata"). La dimensione delle colonne dice qual'è il risultato finale che si produce con questa Operazione, i risultati possibili sono 2: 1) Cessazione assoluta, il mio sistema complesso e aperto a seguito di questa operazione non esiste più. 2) Cessazione Relativa, con questa modalità il sistema complesso e aperto continua ad esistere sotto condizioni differenti. C'è vita dopo una cessazione relativa.
	
	Incrociando queste 2 dimensioni si hanno casistiche
	Cessazione Volontaria, Cessazione Assoluta: LIQUIDAZIONE, L'effetto finale è che a seguito di questa operazione non esiste più alcuna impresa. Tutti gli elementi sono stati venduti singolarmente. Con il ricavato ci pago i debiti e restituisco il capitale ai soci.
	Cessazione Coatta, Cessazione Assoluta: LIQUIDAZIONE (Procedure concorsuali), imposta dalla legge, comunemente si chiama Fallimento. L'azienda non è più in grado di pagare i debiti. Il fallimento provoca la cessazione assoluta, vendita dei singoli beni.
	Cessazione Volontaria, Cessazione Relativa: Ci sono varie cose: "Fusione, Cessione, Scissione, Trasformazione". 
	\begin{itemize}
		\item Trasformazione: Cambio di forma giuridica, operazione mediante la quale una società che ha una certa forma giuridica la cambia e acquista una nuova forma giuridica. Questo perchè la forma giuridica non è data una volta per tutte, così come l'impresa cresce anche la forma giuridica può diventare inadeguata e quindi bisogna cambiarla (Caso CyberPink). Questo viene fatto attraverso la trasformazione. Cambia la forma giuridica ma non i proprietari.
		\item Cessione: Con la vendita una società passa da un Proprietario ad un altro ma continua ad esistere IL valore di Cessione e Liquidazione sono diversi: un sistema complesso è un insieme che ha un valore che è superiore alla somma delle parti, proprio perchè le varie parti hanno una finalità comune che è quello di realizzare un prodotto, venderlo... Questa differenza di valore viene espresso dal termine \textbf{Avviamento}. Capacità che un'impresa ha rispetto ad un altra di produrre profitto. E da che cosa dipende questa capacità in più? Ad esempio: Da dove si trova, Dall'esperienza e dalla capacità dei dipendenti, la storia.
		Chi lo compra, compra una cosa già avviata.
		\item Fusione: Ci sono due forme, Fusione Propria e Incorporazione. 
		\begin{itemize}
			\item Fusione Propria: Immagina avere 3 società. alpha beta e delta. Possono decidere di dar vita a una nuova società mettendo insieme i loro patrimoni. Alpha, Beta e Delta si fondono in Gamma. Con i patrimoni delle prime 3 società.
			\item Fusione per Incorporazione: Si differenzia dalla precedente perchè in questo caso non nasce una nuova società, ma in questo caso beta e delta si fondono in alpha che esisteva e continua ad esistere.
		\end{itemize}
		\item Scorporazione: Crea i così detti spin-off. Nel 2015 Ebay ha fatto uno Spin-Off con PayPal. Immaginiamo una società Z che ha un certo patrimonio, sarà fatto da un insieme di elementi (immobile, sede legale, impanti, crediti...) Questa società cresce molto e crescendo può sviluppare delle linee di Business che diventano autonome (caso di PayPal), creo una nuova società estraendo un patrimonio da quella linea.
		Z continua ad esistere senza più una parte di patrimonio, dando vita ad una o più società
		\item Scissione: Sembra simile allo scorporo ma c'è una differenza importane. Anche in questo caso la partenza è uguale, c'è la società Z che ha un certo patrimonio, con il passare del tempo nascono diverse linee di business e la società può decidere di renderle autonome portandole fuori. In questo caso Z diventa davvero più "magra" (nel caso dello scorporo Z diventa una cassaforte, si prende le azioni). In questo caso le società che escono hanno azioni che possono andare ai soci di Z (Non ha più la cassaforte), sono i soci di Z che devono scegliere se restare soci di Z o scambiare le azioni di Z con le nuove azioni (di Alpha e Beta, che sono le linee di business scisse). Non c'è legame tra Z e Alpha (SCissione parziale, Split-Off). La scissione totale avviene se Z si estingue (Split-Up)
	\end{itemize}
	
	Una Start-Up è un'impresa temporanea, perchè alla ricerca di un business model scalabile e replicabile (E' alla ricerca di un'idea di business che possa avere un mercato, non circoscritto ad una frazione).
	
	Cessazione Coatta, Cessazione Relativa: NON PUO' ESISTERE! 
\end{comment}
	
	\subsection{Operazioni di Natura Straordinaria}
	
	%% NUOVO DEL 25 MARZO
	\textbf{La lezione:} Le competenze ingegneristiche sono il cuore, ma la forma giuridica corretta è il motore che ha permesso a queste aziende di attrarre capitali e crescere.
	
	\subsubsection{Matrice delle Cessazioni}
	Le operazioni straordinarie (non quotidiane) possono essere collocate in una matrice 2x2. 
	\begin{itemize}
		\item \textbf{Righe (Chi decide):} L'operazione può essere su base \textbf{volontaria} o per \textbf{cessazione coatta} (imposta dalla legge, forzata).
		\item \textbf{Colonne (Risultato finale):} 
		\begin{enumerate}
			\item \textbf{Cessazione assoluta:} Il sistema complesso e aperto, a seguito dell'operazione, non esiste più.
			\item \textbf{Cessazione relativa:} Il sistema continua ad esistere sotto condizioni differenti. C'è vita dopo una cessazione relativa.
		\end{enumerate}
	\end{itemize}
	
	\paragraph{Casistiche della matrice:}
	\begin{itemize}
		\item \textbf{Cessazione Volontaria, Cessazione Assoluta: LIQUIDAZIONE.} L'effetto finale è che non esiste più alcuna impresa. Tutti gli elementi sono stati venduti singolarmente. Con il ricavato si pagano i debiti e si restituisce il capitale ai soci.
		\item \textbf{Cessazione Coatta, Cessazione Assoluta: LIQUIDAZIONE (Procedure concorsuali).} Imposta dalla legge, comunemente si chiama \textbf{Fallimento}. L'azienda non è più in grado di pagare i debiti; provoca la vendita dei singoli beni.
		\item \textbf{Cessazione Volontaria, Cessazione Relativa:} Qui rientrano le operazioni di Fusione, Cessione, Scissione e Trasformazione.
		\item \textbf{Cessazione Coatta, Cessazione Relativa: NON PUÒ ESISTERE!}
	\end{itemize}
	
	\subsubsection{Trasformazione, Cessione e Fusione}
	\begin{itemize}
		\item \textbf{Trasformazione:} Cambio di forma giuridica. La società cambia "vestito" perché la vecchia forma è diventata inadeguata alla crescita (Caso CyberPink). Cambia la forma, ma non i proprietari.
		\item \textbf{Cessione:} Con la vendita, la società passa da un Proprietario ad un altro ma continua ad esistere. Il valore di Cessione e quello di Liquidazione sono diversi: un sistema complesso vale più della somma delle parti perché ha una finalità comune. Questa differenza di valore è l'\textbf{Avviamento}. Dipende dalla posizione, dall'esperienza dei dipendenti e dalla storia. Chi compra, compra una cosa già avviata.
		\item \textbf{Fusione:} Si divide in due forme:
		\begin{itemize}
			\item \textbf{Fusione Propria:} Tre società (Alpha, Beta e Delta) decidono di dar vita a una nuova società (Gamma) mettendo insieme i loro patrimoni.
			\item \textbf{Fusione per Incorporazione:} Non nasce una nuova società; Beta e Delta si fondono in Alpha, che esisteva già e continua ad esistere.
		\end{itemize}
	\end{itemize}
	
	\subsubsection{Scorporazione e Scissione (Spin-off, Split-off, Split-up)}
	\begin{itemize}
		\item \textbf{Scorporazione (Spin-off):} Nel 2015 eBay ha fatto uno Spin-Off con PayPal. Immaginiamo una società Z che ha un patrimonio (immobili, impianti, crediti...). Se una linea di business diventa autonoma, creo una nuova società estraendo quel patrimonio. Z continua ad esistere senza quella parte, ma ne possiede le azioni (Z diventa una "cassaforte").
		\item \textbf{Scissione:} Qui la società Z diventa davvero più "magra". A differenza dello scorporo, Z non tiene le azioni della nuova società. Le azioni delle società che escono (Alpha e Beta) vanno direttamente ai \textbf{soci di Z}. Sono i soci che scelgono se restare in Z o scambiare le azioni.
		\begin{itemize}
			\item \textbf{Scissione parziale (Split-Off):} Non c'è legame tra Z e Alpha.
			\item \textbf{Scissione totale (Split-Up):} Avviene se la società Z si estingue del tutto.
		\end{itemize}
	\end{itemize}
	
	\begin{tcolorbox}[colback=blue!5,colframe=blue!75!black,title=Cos'è una Start-up?]
		Una Start-Up è un'impresa \textbf{temporanea}, perché è alla ricerca di un \textit{business model} scalabile e replicabile (un'idea di business che possa avere un mercato non circoscritto ad una frazione).
	\end{tcolorbox}
	
	\subsection{Venture Capital e round di finanziamento}
	
	Il \emph{Venture Capital} (VC) è capitale di rischio per startup ad alto potenziale.%
	Un fondo VC investe denaro in un'impresa in cambio di una quota societaria.%
	Se l'impresa cresce e viene venduta o quotata in borsa (exit), il fondo realizza un profitto elevato; se fallisce, perde l'investimento.
	
	\subsubsection{Il meccanismo}
	
	\begin{enumerate}
		\item Startup con idea promettente ma senza fondi sufficienti per crescere.
		\item Il fondo VC investe capitale (da centinaia di migliaia a milioni di euro).
		\item In cambio: quote societarie, tipicamente tra il 10\% e il 30\%.
		\item L'impresa scala e cresce grazie a capitali, rete e competenze del VC.
		\item Exit: vendita (acquisizione) o IPO (quotazione in borsa); il VC monetizza il suo investimento.
	\end{enumerate}
	
	\subsubsection{Il VC nel Consiglio di Amministrazione}
	
	\begin{itemize}
		\item \textbf{Monitoraggio interno:} controlla direttamente l'andamento e le scelte dell'azienda.
		\item \textbf{Potere di veto:} può bloccare decisioni strategiche contrarie ai suoi interessi.
		\item \textbf{Rete e competenze:} i migliori VC portano contatti, clienti e know-how nel board.
		\item \textbf{Cessione del controllo:} il fondatore perde parte del potere decisionale sull'impresa.
	\end{itemize}
	
	Il VC è accessibile solo con struttura societaria (SRL o SPA), impossibile con l'impresa individuale.
	
	\subsubsection{Round di finanziamento}
	
	\begin{table}[h]
		\centering
		\begin{tabular}{|l|l|l|l|}
			\hline
			\textbf{Round} & \textbf{Fase} & \textbf{Chi investe} & \textbf{Importo tipico} \\
			\hline
			Pre-seed & Idea / prototipo & Fondatori, famiglia, amici & 10k -- 200k € \\
			Seed & MVP, prime validazioni & Business angel, micro-VC & 200k -- 2M € \\
			Serie A & Prodotto funzionante, primi clienti & Fondi VC early-stage & 2M -- 15M € \\
			Serie B & Crescita accelerata, scaling & Fondi VC mid-stage & 15M -- 50M € \\
			Serie C+ & Espansione internazionale & VC late-stage, private equity & 50M+ € \\
			IPO & Quotazione in borsa & Mercato pubblico & --- \\
			\hline
		\end{tabular}
	\end{table}
	
	\paragraph{MVP -- Minimum Viable Product}
	
	È la versione più semplice possibile di un prodotto che permette di testare l'idea sul mercato reale con il minimo sforzo e costo.%
	Concetto chiave: non si costruisce il prodotto perfetto subito, ma lo stretto necessario per rispondere a una domanda fondamentale: «Qualcuno è disposto a pagare per questo?»
	
	«Se non ti vergogni della prima versione del tuo prodotto, l'hai lanciata troppo tardi.» (Reid Hoffman, fondatore di LinkedIn)
	
	\subsection{Oltre il profitto: imprese ibride e sostenibili}
	%% LEZIONE 25 MARZO
	Come vengono distribuiti gli utili? Si va da un mondo for-profit (piena libertà di divisione degli utili), fino ad arrivare ad una not-for-profit in cui non c'è alcuna libertà di divisione degli utili.
	
	Qual'è il fine? Il fine di un'impresa non è il profitto, è uno strumento necessario per far sì che continui ad esistere. 
	
	Come viene distribuito il successo? Da ROI (Indicatore) all'impatto sociale.
	
	\subsubsection{Impresa Forprofit}
	Abbiamo visto le forme giuridiche, SNC, SPA, SAPA. IL successo si misura attraverso dei rating e questi indici vengono calcolati dai dati di bilancio. Responsabilità Sociale: volontaria, dipende dall'impresa, non è imposta per legge (Ferrero SPA, Eni SPA, Luxoticca SPA).
	
	%% FIne
	\subsubsection{Società Benefit}
	
	\begin{itemize}
		\item \textbf{Qualifica giuridica} introdotta dalla Legge 208/2015 (Stabilità 2016).
		\item Obiettivo: profitto + beneficio comune (impatto positivo su persone, comunità, territorio, ambiente).
		\item Obbligo di redigere annualmente una \emph{relazione d'impatto} che descrive le azioni intraprese per il perseguimento del beneficio comune.
		\item Lo scopo di beneficio comune è integrato nello statuto aziendale.
		\item Lo scopo di lucro è previsto insieme allo scopo di beneficio comune: gli utili possono essere distribuiti.
	\end{itemize}
	
	\subsubsection{B-Corp (B-Corporation)}
	Benefit Corporation, non è una forma giuridica italiana. E' una certificazione (Così come ISO 9000).
	\begin{itemize}
		\item \textbf{Certificazione} volontaria rilasciata da B Lab (ente internazionale).
		\item Standard globali di performance sociale e ambientale.
		\item Valutazione rigorosa tramite BIA (B Impact Assessment): punteggio minimo richiesto > 80 punti su 200. Le aree valutate sono la governance, l'ambiente in cui si lavora, la comunità locale in cui la società si trova, il tema ambientale, i clienti...
		\item Focus su trasparenza e responsabilità.
		\item In Italia, le B-Corp certificate sono tenute a modificare il proprio statuto e trasformarsi in Società Benefit per ottenere il rinnovo della certificazione.
	\end{itemize}
	A livello legale non cambia un cazzo, si ha soltanto un certificato. La certificazione vale 3 anni.
	\paragraph{Differenza chiave}
	
	La Società Benefit è una forma legale italiana (status giuridico).%
	La B-Corp è una certificazione internazionale (riconoscimento di performance).%
	Si può essere entrambe: molte aziende italiane sono sia Società Benefit sia B-Corp certificate.
	
	\subsubsection{Social Business}
	%Va scritto, non ho voglia di prendere appunti c'è nelle slide.
	
	\subsubsection{Not For Profit} 
	Chiamato anche Terzo settore (dopo Stato, Mercato).
	%Anche questo
	
	\subsubsection{Impresa Sociale}
	
	\begin{itemize}
		\item Non deve distribuire utili: gli utili devono essere destinati agli scopi statutari o all'incremento del patrimonio.
		\item Scopo principale: utilità sociale, non profitto.
		\item Differenza con Società Benefit: quest'ultima può distribuire utili e persegue \emph{contemporaneamente} profitto e beneficio comune.
	\end{itemize}
	
	\subsubsection{Esempi italiani di successo}
	
	\begin{itemize}
		\item \textbf{Davines Group} (Benefit \& B-Corp): eccellenza nella cosmetica sostenibile; usa ingredienti naturali e investe in progetti di rigenerazione ambientale.
		\item \textbf{Chiesi Farmaceutici} (Benefit \& B-Corp): gruppo internazionale che integra obiettivi di salute dei pazienti con riduzione dell'impatto ambientale.
		\item \textbf{Albergo Etico} (Impresa Sociale): modello di hotel gestito da persone con disabilità; nato ad Asti, oggi è un franchising sociale replicato in Italia.
	\end{itemize}
	
	Messaggio chiave: la sostenibilità non è solo etica, è un vantaggio competitivo che attrae talenti, clienti e investitori.
	
	\subsection{Confronto tra le principali forme giuridiche}
	
	\begin{table}[h]
		\centering
		\small
		\begin{tabular}{|l|l|l|l|l|}
			\hline
			\textbf{Forma} & \textbf{Responsabilità} & \textbf{Capitale} & \textbf{Governance} & \textbf{Ideale per} \\
			\hline
			Impresa Indiv. & Illimitata & Nessuno & Titolare unico & Test / Freelance \\
			\hline
			SNC & Illimitata (Tutti) & Nessuno & Tutti i soci & Piccole attività \\
			\hline
			SAS & Mista & Nessuno & Accomandatari & Soci finanziatori \\
			\hline
			SRL & Limitata & 10.000 € & Flessibile & PMI / Startup \\
			\hline
			SPA & Limitata & 50.000 € & Strutturata & Grandi Imprese \\
			\hline
			SRL Innovativa & Limitata & 1 € (Semplif.) & Flessibile & Startup Tech \\
			\hline
		\end{tabular}
	\end{table}
	
	\subsection{Il ciclo di vita di una startup e le scelte legali}
	
	La forma giuridica non è definitiva: evolve con la crescita dell'impresa.%
	Bisogna sapere quando fare il «salto» alla forma successiva per non frenare la crescita.
	
	\begin{enumerate}
		\item \textbf{Fase 1 -- Idea \& Test:} Impresa Individuale / SRLS (costi minimi, validazione mercato, nessun dipendente).
		\item \textbf{Fase 2 -- Validazione:} SRL / SRL Innovativa (primi dipendenti, investimenti Seed, protezione della proprietà intellettuale).
		\item \textbf{Fase 3 -- Crescita:} SRL Strutturata / SPA (round A/B con VC, espansione, governance complessa).
		\item \textbf{Fase 4 -- Maturità:} SPA (quotazione in borsa -- IPO, mercati globali, azionariato diffuso).
	\end{enumerate}
	
	\subsection{Sintesi: cosa portare a casa}
	
	\begin{enumerate}
		\item \textbf{Scelta strategica:} la forma giuridica non è burocrazia, ma una leva per la crescita e la gestione del rischio.
		\item \textbf{Protezione patrimoniale:} la responsabilità limitata è lo scudo fondamentale per chi fa impresa seriamente.
		\item \textbf{Vantaggi dell'innovazione:} le startup innovative godono di agevolazioni fiscali e burocratiche concrete.
		\item \textbf{Evoluzione continua:} l'assetto societario deve cambiare man mano che l'azienda cresce e matura.
		\item \textbf{Sostenibilità competitiva:} integrare l'impatto sociale (Benefit/B-Corp) è oggi un fattore di successo sul mercato.
	\end{enumerate}
	
	%% Roba nuova.
	\subsection{Le imprese sono Responsabili?}
	%ho perso questa parte **bip**
	
	\begin{comment}\subsection{Caso SatiSpay SpA}
	Fintech (Nasce come una startup nell'ambito finanza e tecnologia) italiana leader nei pagamenti mobile, la sua ambizione è quella di diventare un player europeo, e pone l'obiettivo di crescita esponenziale ed espansione internazionale. Per questo deve raccogliere ingenti capitali (Venture Capital) per marketing e ricerca e sviluppo in un mercato competitivo.
	Partiamo dall'inizio, SatiSpay è una startup e ha questo obiettivo, se dovessimo valutare la forma giuridica più appropriata quale sarebbe la prima domanda che dovremmo porci? (Togliendo il fatto che ora è SpA). Innanzitutto siamo nel mondo for profit, "A che cosa sono interessati i fondi di Venture Capital?" Nel mondo forprofit abbiamo detto che ci sono società di persone e società di capitale (mondo società di persone, responsabilità illimitata. Mondo società di capitali responsabilità limitata). Fondo di venture capital investe con capitale di rischio proprio, quindi preferisce entrare in una forma giuridica con responsabilità limitata. Di conseguenza la scelta del mondo società di capitali è una scelta obbligata. Quale forma potrebbe essere più adatta? Ci sono pro e contro: siccome dobbiamo raccogliere ingenti capitali, la forma per presentarsi meglio sul mercato per questo obiettivo è la SPA (Ci consente anche la quotazione in borsa attraverso la IPO). Qual'è però il rischio di andare verso una SPA? Il rischio è legato alla perdita di controllo (Dalle Azioni, dai Manager). 
	
	\subsection{Caso Exein S.r.l}
	Startup specializzata nella sicurezza del firmware per dispositivi iot. Settore ad alta intensità di Ricerca e Sviluppo. Team composto prevalentemente da ingegneri e ricercatori. E' una SRL Innovativa (Scelta per accedere ad agevolazioni fiscali e flessibilità contrattuale). Per una startup come questa quali sono i vantaggi derivanti dal chiedere la qualifica di startup innovativa? Rivedi 6.6.2.
	Poteva anche prendere un'altra forma giuridica, ma quali sono gli elementi per i quali hanno scelto che la SRL è la forma più adatta? Innanzitutto escludiamo l'SPA perchè ha una complessità tale che non si sposa con il ciclo di vita in cui si trova Extein (E' all'inizio, team composto da ingegneri. Non ha senso assumere una complessità tale non sapendo ancora se possa scalare). Le società di persone sono escluse, siamo in un settore ad alta intensità di ricerca e sviluppo.
	
	\subsection{Caso Davines Group}
	Siamo verso il not for profit. Eccellenza italiana nella cosmetica professionae, azienda familiare con vocazione etica. E ha l'obiettivo di formalizzare l'impegno verso l'ambiente e la società, mettendolo sullo stesso piano del profitto. Che forma giuridica va a cercando? Società Benefit che serve per formalizzare questo impegno nello statuto. Una società benefit può essere anche B-Corp? Sì, non è una forma giuridica, è una certificazione. 
	
	Anche da questi 3 casi emergono i seguenti punti: La forma giuridica è dinamica, deve seguire la strategia aziendale e la vita dell'azienda, ogni scelta ha dei pro e dei contro. Non esiste la forma perfetta. Scegliere una forma giuridica piuttosto che un'altra significa comunicare qualcosa dell'impresa al mercato, è un segnale.
	
	\subsection{Caso Parmalat}
	Parmalat è un colosso, proprietà anche del Parma calcio. Era allora un impero da 7,5 miliardi di euro di fatturato, era presente in 30 paesi nel mondo. A dicembre del 2003 dichiara fallimento, nelle casse di Parmalat mancano 14 Miliardi di euro, 135k piccoli risparmiatori vengono truffati. Come è possibile che nessuno si sia accorto di nulla?.
	Ci poniamo questa domanda: Come un impresa può finanziare le propria crescita? Ci sono due strade principali
	\begin{itemize}
		\item Ricorso all'Equity (Mezzi Propri o Capitale di Rischio). Cedo delle quote di proprietà, non sono io founder ma io più altre persone che vengono sulla barca. Vendo un pezzo dell'impresa, ho un nuovo socio. Condivido i rischi e i profitti sapendo che nessuno di noi ha diritto ad avere indietro alcunché.
		\item Ricorso al Debito (Mezzi di Terzi o Capitale di Credito). Si chiedono dei soldi in prestito, chi mi da i soldi diventa creditore (Banca, Istituto Finanziario, Risparmiatore, Bond...). Dovrò quindi restituire i capitali + gli interessi e manterrò il 100\% della mia impresa. Quanti più debiti faccio, tanto maggiore sarà il vantaggio per i soci.
	\end{itemize}
	
	PS UN bond è un documento che mi dice "Tu mi dai 10, hai diritto ad avere un interesse e il diritto che io ti restituisca 10 alla scadenza".
	
	\subsection{Le Azioni}
	Un'azione è un titolo di credito che rappresenta una quota di proprietà della società. Un'azione attribuisce due diritti: 
	\begin{enumerate}
		\item Diritto di voto: Partecipare alle decisioni e votare
		\item Diritto al dividendo: Ricevere una parte degli utili
	\end{enumerate}
	Non c'è diritto alla restituzione perchè per i soci non esiste!
	Il valore nominale è il numero che è stampato sul foglio. (Così come una banconota da 10 euro, sul foglio c'è scritto 10 euro = Valore Nominale). Un conto è questo, un conto è che cosa io posso comprarci (Potere di Acquisto per la banconota), il valore di mercato è quanto vale in questo momento l'azione. Poi c'è anche il valore di emissione. Cioè il valore che io attribuisco la prima volta che emetto quest'azione sul mercato. La quotazione dell'azione fluttua in funzione di cosa? Il mercato è molto sensibile ai rumors, le aspettative di mercato.
	
	\subsection{Le Obbligazione}
	Sono sempre dei titoli di credito ma danno diritto al rimborso, agli interessi attraverso lo stacco delle cedole e lo status di chi possiede l'obbligazione è quello di creditore (Non c'è alcun diritto di voto, in caso di fallimento gli obbligazionisti vengono prima di tutti gli altri).
	
	\subsection{Diverse tipologie di Azioni}
	\subsubsection{Azioni Ordinarie}
	Azione per definizione, le caratteristiche sono le seguenti
	\begin{itemize}
		\item Diritti amministrativi: Chi possiede queste azioni ha diritto di: partecipare e di votare alle assemblee ordinarie e straordinarie. Chi ha più azioni ordinarie controlla l'azienda. %Le assemblee le scrivo più avanti, sono nelle slide.
		\item Diritti patrimoniali: Diritto al dividendo (solo se c'è utile e l'assemblea decide di distribuirlo) e diritto di opzione in caso di aumento di capitale.
		\item Profilo di Rischio: Massimo, in caso di fallimento gli azionisti ordinari sono gli ultimi ad essere rimborsati (dopo creditori ecc...)
	\end{itemize}
	
	\subsubsection{Azioni di Risparmio}
	Simili alle prima ma hanno più dividendo e meno voce in capitolo
	\begin{itemize}
		\item Nessun diritto di voto: L'azionista di risparmio rinuncia completamente a partecipare alle decisioni. Non vota
		\item Privilegio Patrimoniale: in cambio della rinuncia al voto ottiene un dividendo maggiorato (spesso percentuale minima garantita per legge o statuto, almeno 5\% del valore nominale dell'azione- Devono prendere almeno il 2\% in più delle altre, se questo non fosse vero allora quello che manca per arrivare al 2\% lo recupero il prossimo anno). 
		\item Priorità nel rimborso: in caso di liquidazione, vengono rimborsate prima delle azioni ordinarie.
		
	\end{itemize}
	
	\subsubsection{Azioni Privilegiate}
	E' un compromesso tra le due precedenti
	\begin{itemize}
		\item Voto Limitato: il diritto di voto è limitato soltanto alle sole assmblee straordinarie (Fusioni, modifiche statuarie...). Non votano per il bilancio o il CdA
		\item Privilegio Patrimoniale: In cambio della limitazione del voto, godono di una priorità nella distribuzione degli utili e nel rimborso del capitale rispetto alle ordinarie.
	\end{itemize}
	
	Come possono circolare le azioni privilegiate e ordinarie (che sono diverse da come circolano quelle di risparmio). Nelle azioni ordinarie o privilegiate nel foglio c'è scritto nome e cognome del proprietario (ed è riportato anche nel registro dei soci). Nelle azioni di risparmio non è così (Possesso vale titolo). 
	
	\subsubsection{Azioni a Voto Plurimo e Maggiorato}
	Servono ai Fondatori per mantenere il controllo
	\begin{itemize}
		\item Voto Plurimo: Ad ogni azione corrisponde un voto superiore ad unità, fino a 3 voti per azione. E' uno strumento utilizzato principalmente nelle società non quotate per blindare il controllo
		\item Voto Maggiorato: Premia gli azionisti fedeli, chi detiene l'azione per un periodo continuativo (es. 24 mesi) ottiene 2 voti per azione. Usato nelle società quotate.
	\end{itemize}
	\subsection{Obbligazioni}
	\subsubsection{Obbligazioni Ordinarie}
	E' un contratto di finanziamento in cui chi investe presta del denaro a una impresa o ad uno stato e di conseguenza diventa creditore. Non ha diritto di voto
	Si pone invece il tema della remunerazione, che ha un meccanismo Cedola-Rimborso: Garantisce il pagamento di interessi periodici (cedola fissa) e la restituzione integrale del capitale alla scadenza prestabilita. Il rischio è legato al default dell'emittente e viene misurato dalle agenzie di rating (AAA, BBB)
	
	\subsubsection{Obbligazioni Speciali}
	\begin{itemize}
		\item Obbligazioni Convertibili: Offrono al creditore il diritto (opzione) di trasformare il prestito in azioni della società ad un prezzo prefissato. Si passa da creditori a soci
		\item Obbligazioni Subordinate: In caso di fallimento vengono rimborsate solo dopo aver soddisfatto tutti i creditori ordinari. A fronte del maggior rischio offrono un rendimento più alto.
	\end{itemize}
	\subsection{Caso Parmalat}
	%Da scrivere non ho voglia
	\end{comment}
	
	\subsection{Caso Satispay S.p.A.}
	\textbf{Satispay} è una delle realtà \textit{fintech} (Finance + Technology) più rilevanti in Italia. Nata come startup, la sua evoluzione in S.p.A. è un passaggio strategico obbligato per sostenere un'ambizione di scala europea.
	
	\begin{itemize}
		\item \textbf{Perché la S.p.A.?} La crescita esponenziale richiede capitali ingenti che solo i grandi fondi di \textbf{Venture Capital} o l'eventuale quotazione in borsa (\textbf{IPO}) possono fornire. La forma di S.p.A. offre una maggiore credibilità istituzionale e una struttura di governance che tutela gli investitori internazionali.
		\item \textbf{Il trade-off del controllo:} Per raccogliere centinaia di milioni di euro, i fondatori devono emettere nuove azioni. Questo diluisce la loro quota di proprietà. Il rischio è che i manager e i fondatori perdano il potere decisionale a favore dei grandi fondi d'investimento.
	\end{itemize}
	
	\subsection{Caso Exein S.r.l.}
	Exein opera nel settore della \textit{cybersecurity} per dispositivi IoT (Internet of Things). È un esempio perfetto di come la forma giuridica debba adattarsi alla natura del business.
	
	\begin{itemize}
		\item \textbf{S.r.l. Innovativa:} Hanno scelto questa forma per i vantaggi fiscali legati alla Ricerca e Sviluppo e per la flessibilità nella gestione del team.
		\item \textbf{Perché non S.p.A. subito?} A differenza di Satispay, Exein è in una fase in cui l'agilità è più importante della massa critica di capitali. Una S.p.A. ha costi di gestione e obblighi burocratici (come il Collegio Sindacale sempre obbligatorio) che per una startup \textit{deep-tech} all'inizio sarebbero solo un freno.
	\end{itemize}
	
	\subsection{Caso Davines Group}
	Davines rappresenta il ponte tra il mondo \textit{for-profit} e quello orientato all'impatto sociale.
	
	\begin{tcolorbox}[colback=purple!5,colframe=purple!75!black,title=Focus: Società Benefit vs B-Corp]
		\begin{itemize}
			\item \textbf{Società Benefit (Status Legale):} È una forma giuridica riconosciuta dallo Stato italiano. L'azienda inserisce nel proprio statuto l'obiettivo di generare un impatto positivo sulla società e sull'ambiente, oltre al profitto.
			\item \textbf{B-Corp (Certificazione):} È un "bollino" di qualità rilasciato dall'ente non-profit \textit{B Lab}. Non è una legge, ma un rigido standard internazionale.
		\end{itemize}
		\textit{Nota: Davines è entrambe. La sostenibilità qui non è marketing, ma un asset che attrae i talenti migliori (che oggi cercano uno scopo, non solo uno stipendio).}
	\end{tcolorbox}
	
	\subsection{Il Finanziamento della Crescita: Equity vs Debito}
	Per crescere, l'impresa ha bisogno di "benzina" (capitale). Può ottenerla in due modi:
	
	\begin{enumerate}
		\item \textbf{Equity (Capitale di Rischio):} Si vendono quote della proprietà. 
		\begin{itemize}
			\item \textit{Pro:} Non c'è obbligo di restituzione; i soci rischiano con l'imprenditore.
			\item \textit{Contro:} Si perde il controllo e si dividono gli utili futuri.
		\end{itemize}
		\item \textbf{Debito (Capitale di Credito):} Si chiedono prestiti (banche o obbligazioni).
		\begin{itemize}
			\item \textit{Pro:} L'imprenditore resta padrone al 100\%. Se l'azienda va bene, il rendimento per il socio aumenta (effetto \textit{leva finanziaria}).
			\item \textit{Contro:} Bisogna restituire tutto con gli interessi, anche se l'azienda è in perdita. Troppo debito porta al rischio di \textbf{default}.
		\end{itemize}
	\end{enumerate}
	
	\subsection{Gli Strumenti Finanziari: Azioni e Obbligazioni}
	
	\subsubsection{Le Azioni}
	Rappresentano una quota del capitale sociale. Essere azionista significa essere \textbf{comproprietario}.
	\begin{itemize}
		\item \textbf{Valore Nominale:} Quota di capitale rappresentata (es. 1€).
		\item \textbf{Valore di Mercato:} Prezzo a cui l'azione viene scambiata (dipende dalle aspettative e dai risultati).
		\item \textbf{Diritti:} Voto (amministrativo) e Dividendo (patrimoniale).
	\end{itemize}
	
	\begin{table}[h!]
		\centering
		\small
		\renewcommand{\arraystretch}{1.5}
		\begin{tabular}{|p{3cm}|p{4cm}|p{4cm}|}
			\hline
			\rowcolor{gray!20} \textbf{Tipo di Azione} & \textbf{Diritto di Voto} & \textbf{Privilegio Patrimoniale} \\ \hline
			\textbf{Ordinarie} & Pieno (tutte le assemblee) & Nessuno (dividendo standard) \\ \hline
			\textbf{Privilegiate} & Solo assemblee straordinarie & Priorità nella distribuzione utili \\ \hline
			\textbf{Di Risparmio} & Assente (sono per piccoli risparmiatori) & Dividendo minimo garantito e maggiorato \\ \hline
		\end{tabular}
		\caption{Confronto tipologie di azioni}
	\end{table}
	
	\subsubsection{Le Obbligazioni (Bond)}
	L'investitore non è un socio, ma un \textbf{creditore}. 
	\begin{itemize}
		\item \textbf{Meccanismo:} L'azienda riceve 100 oggi, paga una cedola (interesse) ogni anno e restituisce 100 alla scadenza.
		\item \textbf{Rischio:} Legato alla solidità dell'azienda (Rating). In caso di fallimento, l'obbligazionista viene pagato \textit{prima} dell'azionista.
		\item \textbf{Obbligazioni Convertibili:} Una "opzione" che permette di trasformare il debito in azioni, diventando soci se l'azienda ha successo.
	\end{itemize}
	
	\subsection{Il Caso Parmalat (2003): Il più grande crac finanziario europeo}
	Il caso Parmalat è l'esempio perfetto di cosa succede quando la gestione del debito e la governance falliscono completamente.
	
	\subsubsection{La dinamica del fallimento}
	Parmalat sembrava un impero solido, ma nascondeva un buco di \textbf{14 miliardi di euro}. La strategia fraudolenta si basava su:
	\begin{itemize}
		\item \textbf{Finanziamento tramite Debito Eccessivo:} Per coprire le perdite di società del gruppo (come il Parma Calcio o agenzie di viaggi), l'azienda continuava a emettere \textbf{obbligazioni} sui mercati internazionali.
		\item \textbf{False Comunicazioni Sociali:} Venivano creati falsi conti correnti in paradisi fiscali (es. la società \textit{Bonlat} alle isole Cayman) per dimostrare di avere liquidità che in realtà non esisteva. Celebre il falso documento di Bank of America creato con uno scanner e un po' di colla.
		\item \textbf{Il conflitto di interessi:} Le banche continuavano a vendere obbligazioni Parmalat ai piccoli risparmiatori pur sapendo che l'azienda era sull'orlo del baratro, per rientrare dei propri crediti.
	\end{itemize}
	
	\subsubsection{Cosa ci insegna?}
	\begin{enumerate}
		\item \textbf{Differenza tra Profitto e Cassa:} Parmalat mostrava utili sulla carta, ma non aveva cassa per pagare i debiti.
		\item \textbf{Importanza dei controlli:} Il Collegio Sindacale e le società di revisione fallirono nel loro compito di controllo.
		\item \textbf{Il rischio del debito:} Quando un'impresa usa il debito per coprire perdite operative anziché per investire in crescita, il sistema collassa inevitabilmente (\textit{schema Ponzi} involontario).
	\end{enumerate}
	
	\section{Fonti di finanziamento per le startup}
	\subsection{Business Angels}
	PRivato che investe in startup ad alto potenziale a capitale proprio.
	Interviene nella fase Pre-seed (Normalmente con un importo contenuto tra i 20k e i 500k). Lo fa per investimento personale e passione imprenditoriale.
	
	Siccome interviene nella fase iniziale, svolgono funzioni di mentorin (consigli) e portano oltre il capitale proprio anche una rete di contatti che hanno (utili nelle fasi successive).
	
	Arriva fino ai 3-7 anni
	
	\subsection{Venture Capital}
	Fondo istituzionale. Interviene nelle fasi successive, quando già l'idea (MVP) ha passato il test del mercato. Interviene dal seed alla serie A, B, C e anche oltre (Importi 500k - 50M e oltre). Lo fa per il ritorno finanziario
	
	Ha una struttura molto sviluppata e supporta la fase di scale-up (quando la startup ha avuto semaforo verde e sta uscendo dalla valle della morte).
	
	Prima o poi esco, un exit che varia dai 5-10 anni (Fino all'IPO o all'M\&A)
	\subsection{Corporate Capital}
	Grandi imprese che investono a titolo di mezzi propri in startup. Interviene dalla Serie A, C fino alla fase di crescita. (500K - 50M e oltre) Lo fa perché è interessata all'attività di quella startup (obiettivo strategico).
	
	Porta canali, clienti, tecnologie...
	
	L'exit può arrivare all'acquisizione o all'IPO (Quotazione in borsa).
	
	\section{CyberPeak}
	Ci sono questi soci (Giulia e Marco) iniziano questa società e scelgono come forma giuridica una SNC. Dopodiché c'è una fase di crescita con un venture capital che si trasforma in SRL e poi la fase di maturità (SPA). %Riassunto
	
	\subsection{Prima Fase}
	All'inizio Giulia e Marco perché hanno scelto questa forma giuridica (Società di persone $ \rightarrow $ SNC)?
	Ricordiamo che la SNC ha gestione informale basata sulla fiducia tra i fondatori, in modo da avere decisioni rapide. Non c'è bisogno del capitale minimo.
	
	Ma perchè non una SRL Semplificata? Perchè la SNC è estremamente flessibile, ciascun socio è amministratore. Può scegliere di fare qualcosa senza che l'altro sia d'accordo. 
	
	E qual'è il rischio di questa scelta? Il tema della responsabilità illimitata, in caso di debito devono rispondere con l'intero patrimonio personale.
	
	\subsection{Seconda Fase}
	Entra un fondo di Venture Capital, si passa dalla SNC alla SRL (Richiesto dal fondo di Venture Capital). Ma perchè il fondo ha preteso questa cosa? Il fondo non è disposto a rischiare l'intero suo patrimonio, non è disposto ad entrare come socio in una società di persone (Vogliono responsabilità limitata). Ma perchè proprio SRL? Vuole un posto nel consiglio di amministrazione.
	
	\subsection{Terza Fase}
	La SPA rispetto alla SRL è fortemente regolata (Colleggio sindacale, CdA, aumenta il numero nel CdA). Fase molto critica in cui si sfondando due culture: La flessibilità dei fondatori ma al tempo stesso non hanno soldi per crescere, la cultura finanziaria e strategica di un fondo di venture capital che, sì mette i soldi, ma vuole delle garanzie (anche nei confronti dei terzi).
	
	Due possibilità 1) Il fondo se ne va (exit) ma si rimane senza fondi. 2) Il fondo è un punto di svolta, si cambia o loro non cresceranno mai.
	
	\begin{enumerate}
		
		\item \textbf{Qual è la differenza fondamentale tra un azionista e un obbligazionista?}
		\begin{itemize}
			\item[a)] L'azionista riceve un interesse fisso, l'obbligazionista un dividendo variabile.
			\item[b)] \textbf{L'azionista è un proprietario dell'azienda, l'obbligazionista è un creditore. (CORRETTA)}
			\item[c)] L'azionista ha diritto al rimborso del capitale a scadenza, l'obbligazionista no.
			\item[d)] L'azionista non ha diritto di voto, l'obbligazionista sì.
		\end{itemize}
		\textit{Perché le altre sono errate:} La \textbf{a} inverte i ruoli (l'interesse è per il debito, il dividendo per il capitale). La \textbf{c} sbaglia il soggetto: solo il creditore (obbligazionista) ha un termine di rimborso prestabilito. La \textbf{d} inverte i diritti amministrativi: il voto spetta a chi rischia il capitale proprio (azionista).
		
		\item \textbf{Cosa rappresenta un'azione?}
		\begin{itemize}
			\item[a)] Un prestito fatto all'azienda.
			\item[b)] \textbf{Una quota di proprietà dell'azienda. (CORRETTA)}
			\item[c)] Il diritto a ricevere un interesse fisso.
			\item[d)] Un contratto di lavoro a tempo indeterminato.
		\end{itemize}
		\textit{Perché le altre sono errate:} La \textbf{a} e la \textbf{c} descrivono un'obbligazione. La \textbf{d} confonde un titolo finanziario con un rapporto di lavoro subordinato.
		
		\item \textbf{In caso di fallimento di un'azienda, chi viene rimborsato per primo?}
		\begin{itemize}
			\item[a)] Gli azionisti, perché sono i proprietari.
			\item[b)] I manager, per il lavoro svolto.
			\item[c)] \textbf{Gli obbligazionisti, perché sono creditori. (CORRETTA)}
			\item[d)] Lo Stato, sempre e comunque.
		\end{itemize}
		\textit{Perché le altre sono errate:} La \textbf{a} è l'errore più grave: i proprietari sono gli ultimi a essere pagati (investitori residui). La \textbf{b} è errata perché i manager non hanno priorità assoluta rispetto ai creditori. La \textbf{d} è imprecisa: lo Stato ha privilegi fiscali, ma certi creditori garantiti (es. banche con ipoteca) possono precederlo su certi beni.
		
		\item \textbf{Qual è il principale obiettivo della Corporate Governance?}
		\begin{itemize}
			\item[a)] Massimizzare il profitto a brevissimo termine.
			\item[b)] \textbf{Assicurare che l'azienda sia gestita in modo trasparente, efficiente e responsabile. (CORRETTA)}
			\item[c)] Garantire che solo il fondatore possa prendere decisioni.
			\item[d)] Aumentare il numero di dipendenti dell'azienda.
		\end{itemize}
		\textit{Perché le altre sono errate:} La \textbf{a} è l'antitesi della buona governance, che guarda alla sostenibilità nel tempo. La \textbf{c} descrive un accentramento di potere che la governance invece cerca di bilanciare. La \textbf{d} è un obiettivo occupazionale o di business, non di sistema di controllo e gestione.
		
		\item \textbf{Nel modello di governance "Tradizionale" italiano, quali sono i due organi principali che si occupano di gestione e controllo?}
		\begin{itemize}
			\item[a)] Consiglio di Sorveglianza e Consiglio di Gestione.
			\item[b)] Assemblea dei Soci e Amministratore Delegato.
			\item[c)] \textbf{Consiglio di Amministrazione e Collegio Sindacale. (CORRETTA)}
			\item[d)] Comitato per il Controllo e Consiglio di Amministrazione.
		\end{itemize}
		\textit{Perché le altre sono errate:} La \textbf{a} si riferisce al modello \textit{dualistico}. La \textbf{b} cita l'organo sovrano (Assemblea) e una carica singola, non la coppia di organi di gestione/controllo. La \textbf{d} si riferisce al modello \textit{monistico}.
		
	\end{enumerate}
	\begin{enumerate}[start=6]
		
		\item \textbf{Il caso Parmalat ha dimostrato un colossale fallimento della governance. Quale dei seguenti è stato il più grave conflitto di interessi?}
		\begin{itemize}
			\item[a)] La società di revisione era anche consulente dell'azienda.
			\item[b)] Le banche prestavano soldi pur conoscendo i rischi.
			\item[c)] \textbf{Il fondatore Calisto Tanzi ricopriva i ruoli di azionista di controllo, Presidente e Amministratore Delegato. (CORRETTA)}
			\item[d)] I dipendenti non hanno denunciato le irregolarità.
		\end{itemize}
		\textit{Perché le altre sono errate:} Sebbene la \textbf{a} e la \textbf{b} siano state concause gravissime del crac, la \textbf{c} rappresenta il conflitto di interessi ``originario'' della struttura di governance: l'assenza totale di separazione tra chi decide, chi esegue e chi controlla. La \textbf{d} riguarda l'etica professionale, non un conflitto di interessi strutturale.
		
		\item \textbf{Cosa si intende per modello di governance "Dualistico"?}
		\begin{itemize}
			\item[a)] Un modello in cui le decisioni vengono prese da due Amministratori Delegati.
			\item[b)] \textbf{Un modello, di derivazione tedesca, con una netta separazione tra un Consiglio di Sorveglianza (controllo) e un Consiglio di Gestione (gestione). (CORRETTA)}
			\item[c)] Un modello in cui il controllo è affidato a un comitato interno al Consiglio di Amministrazione.
			\item[d)] Un modello in cui metà del CdA è composto da donne e metà da uomini.
		\end{itemize}
		\textit{Perché le altre sono errate:} La \textbf{a} descrive una semplice co-amministrazione. La \textbf{c} descrive il modello \textit{monistico} (sistema anglosassone). La \textbf{d} confonde la governance societaria con le politiche di diversità di genere (quote rosa).
		
		\item \textbf{Un investitore che cerca un rendimento potenzialmente elevato ed è disposto ad accettare un rischio alto, quale strumento dovrebbe preferire?}
		\begin{itemize}
			\item[a)] Obbligazioni societarie a tasso fisso.
			\item[b)] Buoni del Tesoro Poliennali (BTP).
			\item[c)] \textbf{Azioni di una startup tecnologica. (CORRETTA)}
			\item[d)] Conto deposito bancario.
		\end{itemize}
		\textit{Perché le altre sono errate:} La \textbf{a}, la \textbf{b} e la \textbf{d} sono strumenti a basso o medio rischio con rendimenti limitati. Le azioni di startup (\textbf{c}) offrono il binomio rischio/rendimento più elevato: alta probabilità di perdita totale ma potenziale di crescita esponenziale.
		
		\item \textbf{Quale dei seguenti NON è un diritto tipico di un azionista ordinario?}
		\begin{itemize}
			\item[a)] Diritto di voto nell'assemblea dei soci.
			\item[b)] Diritto a ricevere una parte degli utili (dividendo), se distribuiti.
			\item[c)] \textbf{Diritto alla restituzione garantita del capitale investito. (CORRETTA)}
			\item[d)] Diritto a ispezionare alcuni libri sociali.
		\end{itemize}
		\textit{Perché le altre sono errate:} I punti \textbf{a}, \textbf{b} e \textbf{d} sono i diritti standard previsti dal codice civile. La \textbf{c} è errata perché l'azionista è un investitore di rischio; non esiste alcuna garanzia di restituzione del capitale, a differenza dei creditori.
		
		\item \textbf{Perché il "trucco" di Parmalat, basato sull'emissione continua di obbligazioni, ha potuto funzionare per anni?}
		\begin{itemize}
			\item[a)] Perché le obbligazioni erano garantite dallo Stato italiano.
			\item[b)] Perché gli interessi offerti erano talmente bassi da non destare sospetti.
			\item[c)] \textbf{Perché l'azienda falsificava i bilanci, mostrando una solidità finanziaria inesistente che ingannava i risparmiatori. (CORRETTA)}
			\item[d)] Perché gli investitori erano solo banche e fondi speculativi, non piccoli risparmiatori.
		\end{itemize}
		\textit{Perché le altre sono errate:} La \textbf{a} è falsa (erano titoli privati). La \textbf{b} è errata (i rendimenti erano normali). La \textbf{d} è storicamente falsa, poiché il crac coinvolse migliaia di piccoli risparmiatori che ritenevano Parmalat un investimento sicuro basandosi su bilanci falsi (\textbf{c}).
		
	\end{enumerate}
	
	\section{Fondamenti del Bilancio}
	
	Studiamo i fondamenti del bilancio partendo da un caso reale, quello di De Cecco
	
	\subsection{De Cecco} %Non ho voglia, controllare i bilanci su elearning
	Nel bilancio si dice che ha fatturato (volume di ricavi) di 628M di €. 
	Come mai se l'utile netto è di +11,7M€  nella variazione di cassa si ha un -17M. Se guadagna 12M come mai ha 17M in meno in banca (Da 22.8M a 5.8M)?
	
	Questo modello contabile risponde alla domanda: qual'è lo stato di salute della mia impresa? Ma risponde entro i limiti di validità di questo modello.
	
	Trova infatti ciò che è successo l'anno prima, è uno strumento storico. Risponde a quella domanda utilizzando due variabili
	\begin{itemize}
		\item Reddito: Variabile di flusso. 
		\item Capitale: Variabile di stock. 
	\end{itemize}	
	Metafora del bacino contenente dell'acqua, c'è una diga con le paratie e si crea un bacino in un dato momento con un certo livello dell'acqua in $\phi_0$ pari a $c_0$ (Capitale o "Livello di acqua" che c'è in un determinato momento). Fuori dalla metafora $c_0$ è il capitale della mia impresa fatto dagli elementi che ho (piante, sborra, attrezzatura, crediti, debiti...). 
	Fatto da elementi positivi (attivi, \textbf{attività}) e da elementi passivi (mezzi di terzi, debiti, \textbf{passività}) $= c_0$.
	Si alzano le paratie e fluisce acqua, entra ed esce dal bacino in un tempo $t_0 = 1/1 \quad t_1 = 31/12$ alla fine di $t_1$ chiudiamo e misuriamo l'acqua. L'acqua sarà arrivata ad un livello $c_1$. La variazione di capitale (acqua) è quindi $\varDelta{c}$. I flussi di entrata li chiamo \textbf{Ricavi} che sono entrati da $t_0$ in $t_1$. I flussi in uscita li chiamo \textbf{Costi} che sono usciti da $t_0$ in $t_1$. I ricavi saranno la sommatoria dei ricavi - la sommatoria dei costi (questa sommatoria definisce il Reddito in $t_0$, $t_1$) che potrebbe essere $> 0$ quindi \textbf{Utile} se invece $< 0$ sarà una \textbf{Perdita}.
	
	Nel modello contabile ci sono due rami che si muovono in parallelo:
	Uno che sta davanti a me che vedo si chiama \textbf{Piano finanziario}. Quello che sta dietro si chiama \textbf{Aspetto economico}. 
	
	%PORCODIO MI SONO PERSO DIO CANE SABRINA VAFFANCULO.
	
	
	
	
\end{document}